\exercisegroup

\begin{enumerate}
\def\labelenumi{\arabic{enumi})}
\tightlist
\item
  \begin{enumerate}
  \def\labelenumii{\alph{enumii})}
  \tightlist
  \item
    设 \(\mathbf{E}\) 是一个 Hausdorff 局部凸空间，\(\mathbf{E}'\) 是它的对偶。证明下列性质等价：
  \end{enumerate}
\end{enumerate}

α) E 是次桶型的 (III, 第 44 页, exerc. 7)。

β) E' 中每个强有界子集都是等度连续的。

\(\gamma)\) \(E'\) 中每个强有界子集都是相对弱紧的，并且 \(E\) 的初始拓扑就是 \(\tau(E, E')\) 。

δ) 双对偶 \(E''\) 的强拓扑在 E 上诱导的拓扑与 E 上的初始拓扑相同。

于是，\(\mathbf{E}''\) 的强拓扑的 0 的一个基本邻域系，由对于拓扑 \(\sigma (\mathrm{E}'', \mathrm{E}')\) 而言的闭包组成，这些闭包来自 \(\mathbf{E}\) 的初始拓扑的 0 的一个基本邻域系。

\begin{enumerate}
\def\labelenumi{\alph{enumi})}
\setcounter{enumi}{1}
\tightlist
\item
  证明：若 E 是次桶型的，且它的对偶 \(E'\) 与它的代数对偶 \(E^{*}\) 相同，则 E 的初始拓扑是最细的局部凸拓扑。
\end{enumerate}

¶ 2) a) 证明次桶型空间的任何乘积都是次桶型的。（化归为 Hausdorff 次桶型空间的情形；然后利用 exerc. 1 以及 IV, 第 14 页的命题 15。）

\begin{enumerate}
\def\labelenumi{\alph{enumi})}
\setcounter{enumi}{1}
\tightlist
\item
  给出一个既非有界型亦非桶型的 Hausdorff 次桶型局部凸空间的例子。（按 III, 第 45 页, exerc. 16 的做法，将其中的桶型空间换成次桶型空间，并利用 a)。）
\end{enumerate}

\begin{enumerate}
\def\labelenumi{\arabic{enumi})}
\setcounter{enumi}{2}
\item
  设 E 是一个复 Hausdorff 局部凸空间，\(E_{0}\) 是 E 的底实局部凸空间，并设 \(E^{\prime}\) 与 \(E_{0}^{\prime}\) 分别是 E 与 \(E_{0}\) 的对偶。证明由 \(f \mapsto Rf\) 给出的、从 \(E^{\prime}\) 到 \(E_{0}^{\prime}\) 上的典范 R-线性映射对于 \(E^{\prime}\) 与 \(E_{0}^{\prime}\) 上的 S-拓扑是一个同胚，其中 S 是 E 的有界子集构成的任意一个集合。由此推出从双对偶 \(E^{\prime\prime}\) 到双对偶 \(E_{0}^{\prime\prime}\) 上的典范 R-线性映射的定义，该映射对于弱拓扑 \(\sigma(E^{\prime\prime}, E^{\prime})\) 与 \(\sigma(E_{0}^{\prime\prime}, E_{0}^{\prime})\) ，同样对于强拓扑 \(\beta(E^{\prime\prime}, E^{\prime})\) 与 \(\beta(E_{0}^{\prime\prime}, E_{0}^{\prime})\) ，都是一个同胚；在这个映射下，E（视为嵌入于 \(E^{\prime\prime}\) 中）变换为 \(E_{0}\)（视为嵌入于 \(E_{0}^{\prime\prime}\) 中）。
\item
  设 E 是一个 Hausdorff 次桶型局部凸空间。
\end{enumerate}

\begin{enumerate}
\def\labelenumi{\alph{enumi})}
\tightlist
\item
  证明：若 E 的强对偶 \(E_{b}^{\prime}\) 是有界型的，则 E 的完备化 \(\hat{E}\)（视为 \(E^{\prime*}\) 的一个向量子空间，III, 第 21 页, 定理 2）包含在 E 的双对偶 \(E^{\prime\prime}\) 中。
\item
  我们称 E 是区分的，如果 \(E^{\prime\prime}\) 中每个对于拓扑 \(\sigma(\mathrm{E}^{\prime\prime}, \mathrm{E})\) 有界的子集，都包含在 E 的某个有界子集（对于该拓扑）的闭包之中。证明：E 是区分的，必须且只需它的强对偶 \(E_{b}^{\prime}\) 是桶型的（参见 IV, 第 15 页, 命题 3）。
\end{enumerate}

\begin{enumerate}
\def\labelenumi{\arabic{enumi})}
\setcounter{enumi}{4}
\tightlist
\item
  设 \(\mathbf{E}\) 是一个 Hausdorff 局部凸空间，\(\mathbf{E}'\) 是它的对偶。
\end{enumerate}

\begin{enumerate}
\def\labelenumi{\alph{enumi})}
\item
  要使 \(\mathbf{E}'\) 上的强拓扑与 \(\tau (\mathrm{E}',\mathrm{E})\) 相同，必须且只需 \(\mathbf{E}\) 是半自反的。
\item
  假设 E 是次桶型的。要使 \(E'\) 上的强拓扑与紧收敛拓扑相同，必须且只需 E 是一个 Montel 空间。
\end{enumerate}

¶ 6) 设 F 与 G 是处于分离对偶中的两个向量空间。

\begin{enumerate}
\def\labelenumi{\alph{enumi})}
\tightlist
\item
  证明下列性质等价：
\end{enumerate}

\(\alpha)\) 空间 \(\mathbf{F}\) ，带上一个与 \(\mathbf{F}\) 和 \(\mathbf{G}\) 之间的对偶相容的拓扑，是半自反的；

\(\beta)\) 空间 \(\mathbf{G}\) ，带上 \(\tau (\mathbf{G},\mathbf{F})\) ，是桶型的。

\begin{enumerate}
\def\labelenumi{\alph{enumi})}
\setcounter{enumi}{1}
\tightlist
\item
  证明下列性质等价：
\end{enumerate}

α) 空间 F，带上 τ(F, G)，是自反的；

\(\beta)\) 空间 \(\mathbf{G}\) ，带上 \(\tau (\mathbf{G},\mathbf{F})\) ，是自反的；

γ) F 与 G 分别对于 τ(F, G) 与 τ(G, F) 是桶型的。

\begin{enumerate}
\def\labelenumi{\alph{enumi})}
\setcounter{enumi}{2}
\item
  证明：每个 Hausdorff 局部凸空间 E，带上拓扑 \(\tau(\mathrm{E}, \mathrm{E}')\) ，同构于一个半自反空间 F 关于一个闭子空间的商。（由 III, 第 44 页, exerc. 14, c)，\(E'\) 带上 \(\tau(\mathrm{E}', \mathrm{E})\) 后同构于某个 Hausdorff 桶型空间 G 的一个闭向量子空间 M；取 \(F = G'\) 带上 \(\tau(\mathrm{G}', \mathrm{G})\) ，并利用 IV, 第 10 页的命题 11。）
\item
  设 A 是一个无限集，满足 \(\operatorname{Card}(\mathrm{A}) > \aleph_{1}\) (S, III, §6, exerc. 10)。在乘积空间 \(P = R^{A}\) 中，设 \(E_{0}\) 表示由所有满足如下条件的 \(x = (x_{\alpha})_{\alpha \in \mathrm{A}}\) 组成的处处稠密子空间：除对一个可数的指标集外，均有 \(x_{\alpha} = 0\) ；空间 \(E_{0}\) 是桶型的，并且 P 中由 \(E_{0}\) 与点 \(1_{A}\) （P 中坐标全等于 1 的点）生成的子空间 E 也是桶型的 (III, 第 45 页, exerc. 16)。设 B 是 \(E_{0}\) 中的一个有界子集，它在 P 中的闭包包含 \(1_{A}\) ，并设 J 是 A 的一个基数为 \(\aleph_{1}\) 的子集；设 \(pr_{J}\) 表示从 P 到 \(R^{J}\) 上的投影；证明：B 中存在一个子集 \(B_{J}\) ，其基数 \(\leqslant \aleph_{1}\) ，使得 \(\operatorname{pr}_{J}(B_{J})\) 在 \(R^{J}\) 中的闭包包含 \(1_{J}\) ；于是，集 \(J' \supset J\) 由所有满足如下条件的指标 \(\alpha \in A\) 组成：\(B_{J}\) 中至少有一点，其下标 \(\alpha\) 的坐标 \(\neq 0\) ；该集的基数等于 \(\aleph_{1}\) 。我们定义 \(J_{0} = J\) ，并用归纳法定义 \(J_{n+1} = J'_{n}\) ，令 \(H = \bigcup_{n} J_{n}\) （其基数为 \(\aleph_{1}\) ），并令 \(B_{H} = \bigcup_{n} B_{J_{n}}\) ；证明：\(B_{H}\) 的闭包（从而 B 的闭包）包含点 \((1_{H}, 0) \in R^{H} \times R^{A-H} = P\) ，而该点不属于 \(E_{0}\) 。由此推断：对 E 中每个有界紧集 C，\(C \cap E_{0}\) 在 E 中仍是闭的，并且 E 不是半自反的。
\item
  由 \(d\) 推出：\(\mathrm{E}'\) （即 \(\mathrm{E}\) 的对偶，它可与对偶 \(\mathrm{P}' = \mathbf{R}^{(\mathbf{A})}\) （\(\mathrm{P}\) 的对偶）等同）对于拓扑 \(\tau (\mathrm{E}',\mathrm{E})\) 不是完备的，尽管它对于该拓扑是半自反的（从而是拟完备的）（考虑 \(\mathrm{E}\) 上在 \(\mathrm{E}_0\) 上等于 0 而在点 \(1_{\mathrm{A}}\) 处等于 1 的线性型，并利用 III, 第 21 页, 定理 2）。
\end{enumerate}

\begin{enumerate}
\def\labelenumi{\arabic{enumi})}
\setcounter{enumi}{6}
\item
  设 \((\mathrm{E}_{i})_{i\in\mathbb{I}}\) 是一族 Hausdorff 局部凸空间，P 是诸 \(E_{i}\) 的乘积空间，S 是它们的拓扑直和。证明：要使 P 或 S 是半自反的（相应地，自反的），必须且只需每个 \(E_{i}\) 都是半自反的（相应地，自反的）。
\item
  设 E 是一个 Hausdorff 局部凸空间，它是一个闭向量子空间递增序列 \((E_{n})\) 的严格归纳极限 (II, 第 32 页, 命题 9)。
\end{enumerate}

\begin{enumerate}
\def\labelenumi{\alph{enumi})}
\item
  证明：若每个 \(\mathbf{E}_n\) 的强对偶都是完备的，则 \(\mathbf{E}\) 的强对偶是完备的 (III, 第 20 页, 定理 1)。
\item
  要使 E 是半自反的（相应地，自反的），必须且只需每个 \(E_{n}\) 都是半自反的（相应地，自反的）。
\end{enumerate}

\begin{enumerate}
\def\labelenumi{\arabic{enumi})}
\setcounter{enumi}{8}
\item
  设 \((\mathrm{E}_{\alpha})_{\alpha\in\mathrm{A}}\) 是包含于同一向量空间中的一族 Hausdorff 局部凸空间，它对于关系 \(\supset\) 是有向的，并且当 \(E_{\beta} \subset E_{\alpha}\) 时，\(E_{\beta}\) 的拓扑比 \(E_{\beta}\) 上由 \(E_{\alpha}\) 的拓扑诱导的拓扑更细。设 E 是诸 \(E_{\alpha}\) 的交，带上一个由诸 \(E_{\alpha}\) 的拓扑在 E 上诱导的拓扑取上确界所得的拓扑。证明：若每个 \(E_{\alpha}\) 都是半自反的，则 E 是半自反的（考虑 E 的一个有界子集上的一个超滤子）。
\item
  证明：Montel 空间的任何乘积以及任何拓扑直和都是一个 Montel 空间 \(^{1}\) 。
\item
  设 E 是一个 Hausdorff 局部凸空间，且 E 的强对偶 \(E_{b}^{\prime}\) 是半自反的。a) 证明：在 \(E^{\prime}\) 的每个强有界子集上，由 \(\sigma(E^{\prime}, E^{\prime\prime})\) 与 \(\sigma(E^{\prime}, E)\) 诱导的拓扑相同。
\end{enumerate}

\begin{enumerate}
\def\labelenumi{\alph{enumi})}
\setcounter{enumi}{1}
\tightlist
\item
  由 a) 推出：E 对于拓扑 \(\tau(E, E')\) 是次桶型的（参见 IV, 第 52 页, exerc. 1），并且若 \(\hat{E}\) 是 E 对于该拓扑的完备化，且视为 \(E'^{*}\) 的一个子集 (III, 第 21 页, 定理 2)，则 \(E'' \subset \hat{E}\) 。特别地，若 E 对于 \(\tau(E, E')\) 是拟完备的，则 E 对于该拓扑是自反的。
\end{enumerate}

\(^{1}\) 另一方面，Montel 空间的闭子空间未必是次桶型的，并且 Montel 空间关于闭子空间的商未必是半自反的 (IV, 第 63 页, exerc. 8)。

\begin{enumerate}
\def\labelenumi{\arabic{enumi})}
\setcounter{enumi}{11}
\tightlist
\item
  设 \(\mathbf{E}\) 是一个 Banach 空间，\(\mathbf{E}'\) 是它的对偶。
\end{enumerate}

\begin{enumerate}
\def\labelenumi{\alph{enumi})}
\item
  证明：由 \(x \mapsto d(x, \mathbf{A})\) 给出的、一点 \(x \in \mathbf{E}\) 到一个闭凸集 \(\mathbf{A}\) 的距离，是 \(\mathbf{E}\) 上对于拓扑 \(\sigma(\mathbf{E}, \mathbf{E}')\) 的下半连续函数。
\item
  证明：若 \(\mathrm{E}\) 是自反的，则对每个闭凸子集 \(\mathrm{A}\) （含于 \(\mathrm{E}\) ），存在一点 \(x_0 \in \mathrm{A}\) 使得 \(\| x_0 \|\) 等于 0 到 \(\mathrm{A}\) 的距离（利用 \(a\) ）。若 \(\mathrm{E}\) 的单位球的每个边界点都是极点，则该点是唯一的 (II, 第 54 页, 定义 1)。
\item
  假设 E 是自反的，并设 B 是 E 中一个闭的、凸的且有界的子集；由 a) 与 b) 推出：存在两点 \(x \in \mathbf{A}\) ，\(y \in \mathbf{B}\) ，使得 \(\| x - y \| = d(\mathbf{A}, \mathbf{B})\) （参见 V, 第 71 页, exerc. 8）。
\end{enumerate}

\begin{enumerate}
\def\labelenumi{\arabic{enumi})}
\setcounter{enumi}{12}
\item
  设 \(\mathbf{E}\) 是一个 Banach 空间，\(\mathbf{M}\) 是 \(\mathbf{E}\) 的一个闭向量子空间。证明：若 \(\mathbf{M}\) 与 \(\mathrm{E / M}\) 都是自反的，则 \(\mathbf{E}\) 是自反的。
\item
  设 \(\mathbf{A}\) 是一个无限集。
\end{enumerate}

\begin{enumerate}
\def\labelenumi{\alph{enumi})}
\tightlist
\item
  证明：Banach 空间 \(E = \overline{\mathcal{K}(A)}\) 的强对偶 (IV, 第 47 页, exerc. 1) 可以与 Banach 空间 \(\ell^{1}(A)\) 等同，并且 Banach 空间 \(\ell'(A)\) 的强对偶可以与 Banach 空间 \(\mathcal{B}(A) = \ell^{\infty}(A)\) 等同；由此推出 E 不是自反的，并且 \(E''/E\) 是无限维的（参见 IV, 第 71 页, exerc. 18）。若 A = N，则 E 与 \(E'\) 都是满足第一可数公理的 Banach 空间，但 \(E''\) 不是 (I, 第 25 页, exerc. 1)。
\item
  设 \(B''\) 是 \(E'' = \ell^{\infty}(A)\) 中的单位球，并设 \(B_{0}''\) 是凸集 \(B'' + (B'' \cap E)\) 。证明：\(B_{0}''\) 是 \(E''\) 中对于强拓扑的一个有界闭凸集，其内部非空，但它没有任何极点。若 p 是 \(B_{0}''\) 的度规，证明：带上范数 p 的 \(E''\) 不等距于任何 Banach 空间的对偶，并且 \(B_{0}''\) 对于拓扑 \(\sigma(E'', E')\) 不是闭的（尽管 \(B''\) 对于 \(\sigma(E'', E')\) 是紧的，且 \(B'' \cap E\) 是强闭的）。
\end{enumerate}

¶ 15) 记号与 exerc. 14 中相同。

\begin{enumerate}
\def\labelenumi{\alph{enumi})}
\tightlist
\item
  设 \((x_n')\) 是 \(E'\) 中对于拓扑 \(\sigma(E', E'')\) 收敛于 0 的一个序列；证明：对每个 \(\varepsilon > 0\) ，存在一个有限子集 \(H\) （\(A\) 的），使得 \(\sum_{\alpha \notin H} |x_n'(\alpha)| \leqslant \varepsilon\) 对每个整数 \(n\) 都成立。
\end{enumerate}

（用反证法：若该性质不成立，证明此时存在一个数 \(\delta > 0\) 、整数的一个递增序列 \((n_k)\) 、A 的有限子集的一个递增序列 \((\mathrm{H}_k)\) ，使得 \(\sum_{\alpha \in \mathrm{H}_{k-1}} |x_n'(\alpha)| \leqslant \frac{\delta}{8}\) 对 \(n \geqslant n_k\) 成立，\(\sum_{\alpha \notin \mathrm{H}_k} |x_n'(\alpha)| \leqslant \frac{\delta}{8}\) 对 \(n \leqslant n_k\) 成立，并且

\(\sum_{\alpha \in \mathrm{H}_k - \mathrm{H}_{k-1}} |x_n'(\alpha)| \geqslant \frac{\delta}{2}\) ；证明这导致矛盾（“滑动峰”方法）。）

由此推出序列 \((x_{n}^{\prime})\) 对于强拓扑收敛于 0，尽管强拓扑严格细于拓扑 \(\sigma (\mathrm{E}',\mathrm{E}'')\) 。

\begin{enumerate}
\def\labelenumi{\alph{enumi})}
\setcounter{enumi}{1}
\tightlist
\item
  证明：若 \((x_{n}^{\prime})\) 是 \(E'\) 中对于拓扑 \(\sigma(E', E'')\) 的一个 Cauchy 序列，则它对于该拓扑收敛于 \(E'\) 中的一点；换言之，\(E'\) 对于 \(\sigma(E', E'')\) 是半完备的 (III, 第 7 页) 。（证明：对每个 \(\varepsilon > 0\) ，存在一个有限子集 \(H\) （\(A\) 中的），使得对每个整数 n 都有 \(\sum_{\alpha \notin H} |x_{n}'(\alpha)| \leqslant \varepsilon\) ；按 a) 的方式用反证法，并利用 a)。）
\end{enumerate}

¶ 16) 采用 exerc. 14 的记号，设 E''' 是 \(E'' = \ell^{\infty}(A)\) 的对偶。

\begin{enumerate}
\def\labelenumi{\alph{enumi})}
\item
  设 \(e_{\alpha}\) （\(\alpha \in A\) ）是 \(E''\) 中满足 \(e_{\alpha}(\beta) = \delta_{\alpha \beta}\) （Kronecker 符号）的元素。设 \((\mathbf{K}_n)\) 是 \(A\) 的有限子集的一个序列，两两不相交，并设 \((x_n''')\) 是 \(E'''\) 中的点的一个序列。证明：存在一个严格递增的无穷整数序列 \((n_k)\) （\(>0\) ），使得若令 \(B = \bigcup_k K_{n_k}\) ，则元素 \(y_n' = (x_n'''(\alpha))_{\alpha \in B}\) 属于 \(\ell^1(B)\) 。（设 \(\delta\) 是任意一个 \(>0\) 的数，并设 \((\mathrm{J}_m)_{m \in \mathbb{N}}\) 是 \(\mathbf{N}\) 分成有限集的一个分割。用反证法证明：若 \(x''' \in E'''\) ，则存在一个整数 \(m\) ，使得 \(|\langle x'', x''' \rangle| \leqslant \delta\) 对所有满足下列条件的 \(x'' \in E''\) 成立：\(\| x'' \| \leqslant 1\) 且 \(x''(\alpha) = 0\) ，只要指标 \(\alpha\) 不属于集 \(\bigcup_{n \in J_m} K_n\) 。以适当的方式相继对 \(x_1'', x_2'', ...\) 应用这一结果。）
\item
  由 \(a\) 以及 exerc. 15, a) 推出：若 \((x_{n}^{\prime \prime \prime})\) 是 \(\mathrm{E}'''\) 中对于拓扑 \(\sigma (\mathrm{E}''',\mathrm{E}'')\) 收敛于 0 的一个序列，且 \(\tilde{x}_n^{\prime \prime \prime}\) 是 \(x_{n}^{\prime \prime \prime}\) 在 \(\mathrm{E}''\) 的强闭子空间 E 上的限制，则 \(\lim_{n\to \infty}\| \tilde{x}_n^{\prime \prime \prime}\| = 0\) （在 \(\mathrm{E}'\) 中）。
\item
  由 b) 推出：\(E''\) 的强闭子空间 E 对于强拓扑没有拓扑补。（限于 A = N 的情形，设 \((e_{n}')\) 是 E 上的连续线性型的一个序列，满足 \(\langle x, e_{n}' \rangle = x(n)\) 对所有 \(x \in E\) ；证明序列 \((e_{n}')\) 对于 \(\sigma(\mathrm{E}', \mathrm{E})\) 趋于 0，但 \(e_{n}'\) 不能扩张为连续线性型 \(x_{n}'''\) （\(E''\) 上的），使得序列 \((x_{n}'''')\) 对于 \(\sigma(\mathrm{E}'''', \mathrm{E}'')\) 趋于 0。）
\end{enumerate}

\begin{enumerate}
\def\labelenumi{\arabic{enumi})}
\setcounter{enumi}{16}
\tightlist
\item
  设 E 是一个非自反 Banach 空间，\(E'\) 是它的强对偶，\(E''\) 是 \(E'\) 的强对偶，\(E'''\) 是 \(E''\) 的强对偶，\(E^{IV}\) 是 \(E'''\) 的强对偶。
\end{enumerate}

\begin{enumerate}
\def\labelenumi{\alph{enumi})}
\item
  证明：在 \(E^{\prime\prime}\) 中，\(E^{\prime}\) 与 E 的正交子空间 \(E^{\circ}\) （当 E 视为 \(E^{\prime\prime}\) 的子空间时）是拓扑补，并且 \(E^{\prime\prime}\) 到 \(E^{\prime}\) 上（对于这个分解）的投影是范数为 1 的连续线性映射。与 exerc. 16, c) 相比较，推出 Banach 空间 \(\mathcal{K}(A)\) （作为拓扑向量空间）不同构于任何 Banach 空间的强对偶。
\item
  证明：\(E^{IV}\) 是 \(E^{\circ}\) 与 \(E''\) 的拓扑直和，也是 \(E^{\prime\circ}\) 与 \(E^{\circ\circ}\) 的拓扑直和；我们有 \(E'' \cap E^{\circ\circ} = E\) 。设 v 是 \(E^{IV}\) 到其自身上的线性映射，它在 \(E^{\circ}\) 上是恒同，而在 \(E''\) 上是 \(E''\) 到 \(E^{\circ\circ}\) 上平行于 \(E^{\circ}\) 的投影；证明 v 是一个等距，但对于拓扑 \(\sigma(\mathrm{E}^{\mathrm{IV}}, \mathrm{E}''')\) 不连续。
\end{enumerate}

\begin{enumerate}
\def\labelenumi{\arabic{enumi})}
\setcounter{enumi}{17}
\item
  证明：一个 Banach 空间 E，若其强对偶 \(E'\) 满足第一可数公理，且 E 对于弱化拓扑 \(\sigma(\mathrm{E}, \mathrm{E}')\) 是半完备的 (III, 第 7 页) ，则 E 是自反的（与 IV, 第 54 页, exerc. 15, b) 相比较）。
\item
  设 \(\mathrm{E}\) 是一个 Banach 空间，\(\mathrm{E}'\) 是它的强对偶，\(\mathrm{G}'\) 是 \(\mathrm{E}'\) 的一个对于强拓扑满足第一可数公理的强闭子空间。证明：存在 \(\mathrm{E}\) 的一个可数子集，使得若 \(\mathrm{F}\) 是该子集在 \(\mathrm{E}\) 中生成的闭向量子空间，则 \(\mathrm{G}'\) 等距于 \(\mathrm{F}'\) （\(\mathrm{F}\) 的强对偶）的一个强闭子空间。（设序列 \((x_n')\) 在 \(\mathrm{G}'\) 中强稠密；对每个 \(n\) ，取 \(x_n \in \mathrm{E}\) 使得 \(\| x_n \| \leqslant 1\) 且 \(\langle x_n, x_n' \rangle = \left(1 - \frac{1}{n}\right) \| x_n' \|\) ；证明：强闭子空间 \(F\) 在 \(E\) 中由诸
\end{enumerate}

\[
x _ {n}
\]

生成的强闭子空间 F 即为所求的空间。）

¶ 20) 设 E 是一个 Banach 空间，E' 是它的对偶，B 是 E 中的单位球。要使 B 的每一点对于 B 上由弱化拓扑 σ(E, E') 诱导的拓扑都具有可数的基本邻域系，必须且只需 E' 对于强拓扑满足第一可数公理。（为证明条件是必要的，注意：若 B 的每一点对于弱化拓扑都有可数的基本邻域系，则 B 在 E'' 中对于拓扑 σ(E'', E') 的闭包 B°° 的每一点也是如此。于是存在 E' 中的一个序列 (a'\_n) ，使得 B°° 中 0 的每个 σ(E'', E') 邻域都包含 B°° 与有限多个极集 \{a'\_n\}° 的交；考虑由诸 a'\_n 生成的 E' 的强闭子空间 W' ，以及 W' 在 E'' 中的正交 \(W'^{\circ}\) 。）

\begin{enumerate}
\def\labelenumi{\arabic{enumi})}
\setcounter{enumi}{20}
\tightlist
\item
  设 E 是一个 Banach 空间，\(E'\) 是它的对偶，B 是 E 中的单位球，\(B_{r}'\) 是 \(E'\) 中以 0 为中心、半径为 r 的闭球。
\end{enumerate}

\begin{enumerate}
\def\labelenumi{\alph{enumi})}
\item
  设 \(M_{1}^{\prime}\) ，\(M_{2}^{\prime}\) 是 \(E^{\prime}\) 的两个向量子空间，它们对于弱拓扑 \(\sigma(\mathrm{E}^{\prime}, \mathrm{E})\) 都是处处稠密的。要使 \(\sigma(\mathrm{E}, \mathrm{M}_{1}^{\prime})\) 与 \(\sigma(\mathrm{E}, \mathrm{M}_{2}^{\prime})\) 在 B 上诱导的拓扑相同，必须且只需 \(M_{1}^{\prime}\) 与 \(M_{2}^{\prime}\) 在 \(E^{\prime}\) 中的强闭包相同。
\item
  设 \(M'\) 是 \(E'\) 的一个对于 \(\sigma(E', E)\) 处处稠密的向量子空间；用 \(\mathbf{M}'^{(1)}\) 表示 \(M' \cap B_{1}'\) 在 \(E'\) 中对于拓扑 \(\sigma(E', E)\) 的闭包所生成的向量子空间。要使 \(\mathbf{M}'^{(1)} = E'\) ，必须且只需 \(M' \cap B_{1}'\) 的弱闭包包含一个球 \(B_{r}'\) ，其中 r \textgreater{} 0（利用 \(E'\) 对于强拓扑是桶型的这一事实）。
\item
  设 \(r\) 是这样的数 \(t\) 的上确界：\(\mathbf{M}' \cap \mathbf{B}_1'\) 的弱闭包包含一个球 \(\mathbf{B}_t'\) ；称数 \(r\) 为 \(\mathbf{M}'\) 的特征数。证明：\(r\) 是诸数 \(\sup_{x' \in \mathbf{M}' \cap \mathbf{B}_1'} |\langle x, x' \rangle| / \| x\|\) 的下确界，其中 \(x\) 取遍所有 \(\neq 0\) 且含于 \(E\) 的点（利用 Hahn-Banach 定理）。
\item
  证明：\(1 / r\) 是 \(\| x\|\) 的上确界，其中 \(x\) 取遍 B 在 E 中对于拓扑 \(\sigma (\mathrm{E},\mathbf{M}^{\prime})\) 的闭包（利用 \(c\) ) 与 Hahn-Banach 定理）。
\item
  设 \(M^{\prime\circ}\) 是 \(M^{\prime}\) 在 \(E^{\prime\prime}\) 中的正交；证明：\(r = \inf(\|x + z^{\prime\prime}\|/\|x\|)\) ，其中 \(z^{\prime\prime}\) 取遍 \(M^{\prime\circ}\) ，x 取遍 E 中非零的点（利用 c) 与 Hahn-Banach 定理）。由此推出：要使 \(\mathbf{M}^{\prime(1)} = \mathbf{E}^{\prime}\) ，必须且只需 \(E + M^{\prime\circ}\) 在 \(E^{\prime\prime}\) 中是强闭的（利用 I, 第 17 页的 Banach 定理）。
\end{enumerate}

\(f\) ) 设 \(\mathbf{A} = \mathbf{N}\times \mathbf{N}\) ，\(\mathrm{E} = \overline{\mathcal{K}(\mathbf{A})}\) （参见 IV, 第 54 页, exerc. 14）。在空间 \(\mathrm{E}'' = \ell^{\infty}(\mathrm{A})\) 中，设 \(\mathrm{P}\) 是由所有这样的点 \(x = (x_{ij})\) 组成的向量子空间：\(x_{ij} = x_{0j} / (j + 1)\) 对所有 \(i\geqslant 0\) 。证明：\(\mathbf{P} = \mathbf{M}'^{\circ}\) ，其中 \(\mathbf{M}'\) 是 \(\mathrm{E}'\) 的一个（对于 \(\sigma (\mathrm{E}',\mathrm{E}))\) 处处稠密的向量子空间，但 \(\mathrm{E} + \mathrm{M}'^{\circ} = \mathrm{E} + \mathrm{P}\) 在 \(\mathrm{E}''\) 中不是强闭的；由此推出 \(\mathbf{M}'\) 的特征数为 0。

\begin{enumerate}
\def\labelenumi{\arabic{enumi})}
\setcounter{enumi}{21}
\tightlist
\item
  设 E 是一个 Banach 空间，\(E'\) 是它的对偶，\(M'\) 是 \(E'\) 中的一个强闭子空间，它对于 \(E'\) 上的弱拓扑是处处稠密的。我们称 M' 是不可约的，如果不存在向量子空间 \(N' \neq M'\) （\(M'\) 的），它在 \(E'\) 中是强闭且弱处处稠密的。a) 证明：\(M'\) 不可约的充要条件是正交 \(M'^{\circ}\) （\(M'\) 在 \(E''\) 中的）是 E 的拓扑补（对于 \(E''\) 的强拓扑）。由此推出此时 \(\mathbf{M}^{(1)} = \mathbf{E}'\) (exerc. 21)，并且 E 同构于空间 \(M'\) （带上由 \(E'\) 的强拓扑诱导的拓扑）的强对偶。
\end{enumerate}

\begin{enumerate}
\def\labelenumi{\alph{enumi})}
\setcounter{enumi}{1}
\item
  证明：\(\mathbf{M}'\) 不可约的充要条件是 \(\mathbf{E}\) 中的单位球对于拓扑 \(\sigma (\mathrm{E},\mathbf{M}')\) 是相对紧的（利用 exerc. 21, a)）。
\item
  要使 E （作为拓扑向量空间）同构于某个 Banach 空间的强对偶，必须且只需 \(E'\) 中存在一个不可约子空间。由此给出下述事实的一个新证明：Banach 空间 \(\overline{\mathcal{K}(\mathbf{N})}\) 不同构于任何 Banach 空间的强对偶（参见 IV, 第 55 页, exerc. 16, c)）。
\end{enumerate}

\begin{enumerate}
\def\labelenumi{\arabic{enumi})}
\setcounter{enumi}{22}
\tightlist
\item
  采用与 exerc. 22 相同的记号，假设 \(\mathbf{M}'\) 是不可约的。
\end{enumerate}

\begin{enumerate}
\def\labelenumi{\alph{enumi})}
\item
  要使从 E 到 \(E^{\prime\prime}/M'^{\circ}\) 中的典范映射（它是典范映射 \(E^{\prime\prime} \rightarrow E^{\prime\prime}/M'^{\circ}\) 的限制）是一个 Banach 空间等距，必须且只需 \(M^{\prime}\) 的特征数 (IV, 第 55 页, exerc. 21) 等于 1。这时，称 \(M^{\prime}\) （带上由 \(E^{\prime}\) 的范数诱导的范数）为 E 的预对偶，并且 E 可以（连同其范数）典范地等同于 Banach 空间 \(M^{\prime}\) 的对偶。
\item
  对 E 的每个对于 \(\sigma (\mathrm{E},\mathbf{M}^{\prime})\) 为闭的向量子空间 F ，证明：\(\mathbf{M}'\) 在 F 的对偶 \(\mathrm{F}'\) （它等同于 \(\mathrm{E}^{\prime} / \mathrm{F}^{\circ}\) ）中的典范像是 F 的一个预对偶。
\end{enumerate}

\begin{enumerate}
\def\labelenumi{\arabic{enumi})}
\setcounter{enumi}{23}
\tightlist
\item
  \begin{enumerate}
  \def\labelenumii{\alph{enumii})}
  \tightlist
  \item
    设 \((a_{k})_{1\leqslant k\leqslant n}\) 是赋范空间 \(E\) 中点的一个有限序列，并设 \((\lambda_k)_{1\leqslant k\leqslant n}\) 是数的一个有限序列，诸数 \(>0\) 且满足 \(\sum_{k = 1}^{n}\lambda_{k} < 1\) ；令 \(\mu_{k} = 1 - \sum_{j = 1}^{k - 1}\lambda_{j}\) （对每个 \(k\) ）；那么
  \end{enumerate}
\end{enumerate}

\[
\left\| \sum_ {k = 1} ^ {n} \lambda_ {k} a _ {k} \right\| \leqslant \frac {\lambda_ {n}}{\mu_ {n - 1}} \left\| \mu_ {n - 1} a _ {n} + \sum_ {k = 1} ^ {n - 1} \lambda_ {k} a _ {k} \right\| + \frac {\mu_ {n}}{\mu_ {n - 1}} \left\| \sum_ {k = 1} ^ {n - 1} \lambda_ {k} a _ {k} \right\|.
\]

\begin{enumerate}
\def\labelenumi{\alph{enumi})}
\setcounter{enumi}{1}
\tightlist
\item
  设 \((a_{n})\) 是 E 的单位球中点的一个无穷序列，并设 \((\lambda_{n})\) 是数的一个无穷序列，诸数 \(>0\) 且满足 \(\sum_{n} \lambda_{n} = 1\) 。对每个 \(n > 0\) ，令 \(\mu_{n} = 1 - \sum_{k=1}^{n-1} \lambda_{k}\) ，并令 \(b_{n} = \sum_{k=1}^{n-1} \lambda_{k} a_{k} + \mu_{n} a_{n}\) ；证明：对所有 \(n \geqslant 1\) ，我们有
\end{enumerate}

\[
\left\| \sum_ {k = 1} ^ {n} \lambda_ {k} a _ {k} \right\| \leqslant \mu_ {n} \sum_ {k = 1} ^ {n} \frac {\lambda_ {k} \| b _ {k} \|}{\mu_ {k - 1} \mu_ {k}}.
\]

（用归纳法应用 a)。）

\begin{enumerate}
\def\labelenumi{\alph{enumi})}
\setcounter{enumi}{2}
\tightlist
\item
  设 \((\mathrm{C}_n)\) 是 E 中凸集的一个递减序列，包含于单位球中，且满足 \(d(0,\mathrm{C}_1)\geqslant \theta >0\) （从而更有 \(d(0,\mathrm{C}_n)\geqslant \theta\) 对所有 \(n\) ）。设 \((\lambda_{n})\) 是数的一个序列，诸数 \(>0\) 且满足 \(\sum_{n}\lambda_{n} = 1\) 。证明：存在一个数 \(\alpha\) 使得 \(\theta \leqslant \alpha \leqslant 1\) ，以及 E 的点的一个序列 \((x_{n})\) ，使得 \(x_{n}\in \mathbf{C}_{n}\) 对所有 \(n\) ，\(\| \sum_{n}\lambda_{n}x_{n}\| = \alpha\) ，且对所有 \(n\)
\end{enumerate}

\[
\left\| \sum_ {k = 1} ^ {n} \lambda_ {k} x _ {k} \right\| \leqslant \alpha (1 - \theta \sum_ {j = n + 1} ^ {\infty} \lambda_ {j}).
\]

（取 \(x_{1}\) 使 \(\| x_{1}\|\) 任意接近 \(d(0, \mathrm{C}_{1})\) ，然后用归纳法取 \(x_{n}\) ，使 \(\left\| \sum_{k=1}^{n-1} \lambda_{k} x_{k} + \left(\sum_{j=n}^{\infty} \lambda_{j}\right) x_{n}\right\|\) 任意接近诸数 \(\left\| \sum_{k=1}^{n-1} \lambda_{k} x_{k} + \left(\sum_{j=n}^{\infty} \lambda_{j}\right) y\right\|\) 的下确界，其中 \(y\) 取遍 \(\mathrm{C}_n\) 。然后利用 \(b)\) 。）

¶ 25) 设 E 是一个满足第一可数公理的 Banach 空间。证明下列性质等价：

α) E 不是自反的。

\(\beta)\) 对每个数 \(\theta\) （满足 \(0 < \theta < 1\) ），存在一个序列 \((x_n')\) （在 \(\mathbf{E}'\) 中），使得 \(\| x_n'\| \leqslant 1\) 对所有 \(n\) ，序列 \((x_n')\) 对于 \(\sigma (\mathrm{E}',\mathrm{E})\) 收敛于 0，并且 0 到由诸 \(x_{n}^{\prime}\) 生成的凸集的距离 \(\geqslant \theta\) 。

\(\gamma)\) 对每个数 \(\theta\) （满足 \(0 < \theta < 1\) ）以及每个序列 \(\lambda_{n}\) （诸数 \(>0\) ，满足 \(\sum_{n} \lambda_{n} = 1\) ），存在一个数 \(\alpha\) 使得 \(\theta \leqslant \alpha \leqslant 1\) ，以及点的一个序列 \((y_{n}')\) ，其点属于

E' ，使得 \(\| y_n'\| \leqslant 1\) 对所有 \(n\) ，\(\left\| \sum_{n}\lambda_ny_n'\right\| = \alpha\) ，且 \(\left\| \sum_{k=1}^{n}\lambda_ky_k'\right\| \leqslant \alpha (1 - \theta \sum_{j=n+1}^{\infty}\lambda_j)\) 对所有 \(n\) 。

\(\delta)\) 存在 \(z' \in \mathrm{E}'\) ，使得没有任何 \(x \in \mathrm{E}\) 满足 \(\left|\langle x, z' \rangle\right| = \| x \| \cdot \| z'\|\) （James-Klee 定理）。

（要看 \(\alpha\) ) 蕴含 \(\beta\) )，注意存在 \(z'' \in \mathrm{E}''\) 使得 \(\| z'' \| < 1\) 且 \(d(z'', \mathrm{E}) > \theta\) 。若 \((x_n)\) 是 \(\mathrm{E}\) 中的一个处处稠密序列，求一个序列 \((x_n')\) （在 \(\mathrm{E}'\) 中），使得 \(\| x_n' \| < 1\) ，且 \(\langle x_k, x_n' \rangle = 0\) （对 \(k \leqslant n\) ）以及 \(\langle x_n', z'' \rangle = \theta\) 。要看 \(\beta\) ) 蕴含 \(\gamma\) )，利用 exerc. 24。要看 \(\gamma\) ) 蕴含 \(\delta\) )，证明对每个 \(x \in \mathrm{E}\) 有 \(|\sum_{n} \lambda_n \langle x, y_n' \rangle| < \alpha\) （利用 \(\gamma\) ）。）

\begin{enumerate}
\def\labelenumi{\arabic{enumi})}
\setcounter{enumi}{25}
\tightlist
\item
  我们称局部凸空间 E 具有性质 (GDF)，如果从 E 到某个 Banach 空间 F 中的每个满足下列性质的线性映射 u 都是连续的：在乘积空间 \(E \times F\) 中，u 的图像 \(\Gamma\) 中收敛点列的每个极限仍属于 \(\Gamma\) 。每个 Fréchet 空间都具有性质 (GDF) (I, 第 19 页, 推论 5)；Fréchet 空间族的每个归纳极限也是如此 (II, 第 34 页, 命题 10)。证明：每个具有性质 (GDF) 的 Hausdorff 局部凸空间都是桶型的。（设 V 是 E 中的一个桶，q 是它的度规，H 是 E 带上这个半范数所对应的 Hausdorff 空间；证明典范映射 \(\pi\) （从 E 到完备化 \(\hat{H}\) 中的）是连续的，这要利用性质 (GDF) 以及每个线性型 \(x' \in V^{0}\) 都可以唯一地扩张成 \(\hat{H}\) 上的连续线性型这一事实（这些线性型所成的集是 \(\hat{H}\) 的对偶中的单位球）。）
\end{enumerate}

\exercisegroup

¶ 1) 设 E 是一个可度量的局部凸空间，\(E_{b}^{\prime}\) 是它的强对偶。若 \(E_{b}^{\prime}\) 是可度量的，证明 E 的拓扑可以由单独一个范数定义（利用 III, 第 37 页, exerc. 2 与第 38 页, exerc. 5，以及 E 是有界型的这一事实）。

\begin{enumerate}
\def\labelenumi{\arabic{enumi})}
\setcounter{enumi}{1}
\tightlist
\item
  次桶型空间是半桶型的。我们称一个局部凸空间为 (DF) 空间，如果它是半桶型的，并且其典范有界结构 (III, 第 3 页, 定义 5) 具有可数基。每个赋范空间以及赋范空间序列的每个严格归纳极限 (II, 第 33 页) 都是 (DF) 空间。Fréchet 空间的每个强对偶都是 (DF) 空间。
\end{enumerate}

\begin{enumerate}
\def\labelenumi{\alph{enumi})}
\item
  (DF) 空间的强对偶是 Fréchet 空间。
\item
  设 \(\mathrm{E}\) 是一个 (DF) 空间，\((\mathbf{A}_n)\) 是 \(\mathrm{E}\) 中有界、凸、均衡且闭的子集的一个递增序列，使得 \(\mathrm{E}\) 的每个有界子集都被诸 \(\mathbf{A}_n\) 之一吸收。
\end{enumerate}

  设 U 是诸 \(A_{n}\) 的并；证明：U 在 E 中的闭包 \(\overline{U}\) 恰好是所有这样的 \(x \in E\) 组成的集：\(\lambda x \in U\) 对 \(0 \leqslant \lambda < 1\) 。（若 \(x \notin \lambda U\) 对某个 \(\lambda > 1\) ，则对每个 n，存在线性型 \(x_{n}' \in E'\) 使得 \(x_{n}' \in A_{n}^{\circ}\) 且 \(\langle x, x_{n}' \rangle = \lambda\) ，而序列 \((x_{n}')\) 是等度连续的，从而有一个弱极限点。）

\begin{enumerate}
\def\labelenumi{\alph{enumi})}
\setcounter{enumi}{2}
\tightlist
\item
  证明：若一个 (DF) 空间是桶型的，则它也是有界型的（参见 III, 第 44 页, exerc. 13, b)）。给出非超有界型但是有界型且桶型的 (DF) 空间的例子 (III, 第 46 页, exerc. 22)，以及非桶型但是有界型的 (DF) 空间的例子。
\end{enumerate}

\begin{enumerate}
\def\labelenumi{\arabic{enumi})}
\setcounter{enumi}{2}
\tightlist
\item
  设 \(E\) 是一个可度量的局部凸空间，\(E_b'\) 是它的强对偶。
\end{enumerate}

\begin{enumerate}
\def\labelenumi{\alph{enumi})}
\item
  证明：每个凸均衡子集 \(V'\) （属于 \(E_{b}'\) ）只要吸收 \(E_{b}'\) 的强有界子集，就包含一个（对于强拓扑的）桶，该桶吸收 \(E_{b}'\) 的强有界子集。（设 \((\mathbf{K}_{n}')\) 是 \(E_{b}'\) 的典范有界结构的一个可数基，并设 \(\lambda_{n}\) 使得 \(\lambda_{n}K_{n}' \subset \frac{1}{2}V'\) ；把 exerc. 2, b) 应用于序列 \(A_{n}'\) ，其中 \(A_{n}'\) 是诸 \(\lambda_{j}K_{j}'\) （\(j \leqslant n\) ）的并的凸包。）
\item
  由 \(a\) ) 推出下列性质等价：
\end{enumerate}

α) E 是区分的 (IV, 第 52 页, exerc. 4)。

β) \(E_{b}^{\prime}\) 是次桶型的 (III, 第 44 页, exerc. 7)。

\(\gamma)\) \(\mathrm{E}_b^{\prime}\) 是有界型的。

δ) \(E_{b}^{\prime}\) 是桶型的。

ε) \(E_{b}^{\prime}\) 是超有界型的 (III, 第 45 页, exerc. 19)。

\begin{enumerate}
\def\labelenumi{\alph{enumi})}
\setcounter{enumi}{2}
\tightlist
\item
  证明：若 \(E_{b}^{\prime}\) 是自反的，则 \(\hat{E} = E''\) （它显然是自反的）（参见 IV, 第 52 页, exerc. 4 及第 53 页, exerc. 11）。
\end{enumerate}

\begin{enumerate}
\def\labelenumi{\arabic{enumi})}
\setcounter{enumi}{3}
\tightlist
\item
  设 E 是一个 Hausdorff 局部凸空间，\(E'\) 是它的对偶。若 M 是 E 的一个可度量且区分的 (IV, 第 52 页, exerc. 4) 闭向量子空间，则强拓扑 \(\beta(\mathrm{E}'/\mathrm{M}^{\circ}, \mathrm{M})\) 是按 \(M^{\circ}\) 从强拓扑 \(\beta(\mathrm{E}', \mathrm{E})\) 得到的商拓扑（利用 exerc. 3, b) 与 IV, 第 51 页, exerc. 22, b)）。
\end{enumerate}

¶ 5) 对每个整数 n \textgreater{} 0，设 \(a^{(n)}\) 是二重序列 \((a_{pq}^{(n)})\) （ \(p \in N\) ，\(q \in N\) ），满足：\(a_{pq}^{(n)} = q\) 若 \(p \leqslant n\) ，且 \(a_{pq}^{(n)} = 1\) 若 p \textgreater{} n。~设 E 是由所有满足如下条件的实数二重序列 \(x = (x_{pq})_{(p,q) \in \mathbf{N} \times \mathbf{N}}\) 组成的向量空间：对每个整数 n \textgreater{} 0，数 \(r_n(x) = \sum_{p,q} a_{pq}^{(n)} |x_{pq}|\)

是有限的。若给 E 赋予由诸半范数 \(r_n\) 定义的拓扑，则 E 是一个满足第一可数公理的 Fréchet 空间 (IV, 第 47 页, exerc. 1, c))；E 的对偶 E' 可以与由所有满足如下条件的实数二重序列 \(x' = (x_{pq}')\) 组成的空间等同：至少对一个指标 \(n\) ，存在 \(k_n > 0\) 使得 \(|x_{pq}'| \leqslant k_n a_{pq}^{(n)}\) 对每一对 \((p, q)\) ；并且 \(\langle x, x' \rangle = \sum_{p, q} x_{pq} x_{pq}'\) 。

(IV, 第 47 页, exerc. 1, c))。

对每个整数 \(p_0 > 0\) 以及整数的每个序列 \((m_p)\) （诸整数 \(> 0\) ），设 \(\mathrm{J}(p_0; (m_p))\) 是由所有整数对 \(p > 0\) ，\(q > 0\) （满足 \(p \geqslant p_0\) 且 \(q \geqslant m_p\) ）组成的集；设 \(\mathfrak{V}\) 是 \(\mathbf{N} \times \mathbf{N}\) 上由诸集 \(\mathrm{J}(p_0; (m_p))\) 组成的滤子基，并设 \(\mathfrak{F}\) 是一个比以 \(\mathfrak{V}\) 为基的滤子更细的超滤子。

\begin{enumerate}
\def\labelenumi{\alph{enumi})}
\item
  证明：对所有 \(x' = (x_{pq}') \in \mathrm{E}'\) ，二重序列 \((x_{pq}')\) 有一个极限 \(u(x')\) （关于超滤子 \(\mathfrak{F}\) ）；若 \(\mathrm{V}_n\) 是 \(\mathrm{E}\) 中由 \(r_n(x) \leqslant 1\) 定义的 0 的邻域，则 \(|u(x')| \leqslant 1\) 对所有 \(x' \in \mathrm{V}_n^\circ\) 。
\item
  设 \(U'\) 是 \(E'\) 中对于强拓扑的 0 的一个邻域，它是凸的、均衡的且弱闭的，并对每个 n，设 \(\alpha_{n} > 0\) 使得 \(\alpha_{n} V_{n}^{\circ} \subset U'\) 。对每个整数 p \textgreater{} 0，设 \(m_{p}\) 是使得 \(2^{p+1} \leqslant \alpha_{p} m_{p}\) 的整数，并设 \(x' = (x'_{pq})\) 是满足如下条件的二重序列：\(x'_{pq} = 0\) 对 q \textless{} \(m_{p}\) ，\(x'_{pq} = 2\) 对 \(q \geqslant m_{p}\) 。证明：\(x' \in U'\) 但 \(u(x') = 2\) ；由此推出 u 在 \(E'\) 中不是强连续的，尽管它在 \(E'\) 的每个有界子集上是有界的。推断 (IV, 第 58 页, exerc. 3)：E 不是区分的，从而强对偶 \(E'_{b}\) 是一个非次桶型的 (DF) 空间。
\item
  利用 b) 构造一个 Fréchet 空间 F 的闭子空间 M 的例子，使得强拓扑 \(\beta(\mathrm{F}^{\prime}/\mathrm{M}^{\circ}, \mathrm{M})\) 不同于按 \(M^{\circ}\) 从强拓扑 \(\beta(\mathrm{F}^{\prime}, \mathrm{F})\) 得到的商拓扑（把 E 嵌入到一个 Banach 空间的可数乘积中）。
\end{enumerate}

¶ 6) a) 设 E 是一个 (DF) 空间 (IV, 第 57 页, exerc. 2)，并设 U 是一个凸集，使得对 E 的每个有界、凸且均衡的子集 A，U ∩ A 都是 0 的一个邻域（对于 E 的拓扑在 A 上诱导的拓扑）。证明：U 是 A 中 0 的一个邻域。（设 (A \(_{n}\) ) 是 E 的典范有界结构的一个可数基 (III, 第 3 页, 定义 5)。证明：可以用归纳法定义一个序列 ( \(\lambda_{n}\) ) （由数 \textgreater{} 0 组成）以及一个序列 (V \(_{n}\) ) （由闭的、凸且均衡的

0 的邻域组成），使得 \(\lambda_{n}\mathrm{A}_{n} \subset \frac{1}{3}\mathrm{U}\) ，\(\lambda_{n}\mathrm{A}_{n} \subset \bigcap_{j=1}^{\infty} \mathrm{V}_{j}\) ，\(\mathrm{V}_{n} \cap \mathrm{A}_{n} \subset \mathrm{U}\) 对每个 \(n\) 。

首先证明：若 \(\lambda_{j}\) 与 \(\mathrm{V}_j\) 已对 \(j\leqslant n\) 构造出来，则可找到 \(\lambda_{n + 1}\) 使得 \(\lambda_{n + 1}\mathrm{A}_{n + 1}\subset \frac{1}{3}\mathrm{U}\) 且 \(\lambda_{n + 1}\mathrm{A}_{n + 1}\subset \mathrm{V}_j\) 对 \(j\leqslant n\) 。其次证明：可找到 \(\mathrm{V}_{n + 1}\) 使得 \(\lambda_j\mathrm{A}_j\subset \mathrm{V}_{n + 1}\) 对 \(j\leqslant n + 1\) 且 \(\mathrm{V}_{n + 1}\cap \mathrm{A}_{n + 1}\subset \mathrm{U}\) ；为此，设 A 表示诸 \(\lambda_{j}\mathrm{A}_{j}\) （\(j\leqslant n + 1\) ）的凸包，证明可取 \(\mathrm{V}_{n + 1} = \overline{\mathrm{A} + \mathrm{V}}\) ，其中 V 是一个适当的凸均衡 0 邻域；我们指出，为此只需证明：若 \(\mathbf{B} = \mathbf{A}_{n + 1}\cap \mathbb{C}\) U，则 0 不在集 \(\mathbf{B} + 2\mathbf{A}\) 的闭包中。

\begin{enumerate}
\def\labelenumi{\alph{enumi})}
\setcounter{enumi}{1}
\tightlist
\item
  由 a) 推出：若 u 是从 E 到某个局部凸空间 F 中的线性映射，且 u 在 E 的每个有界子集上的限制都是连续的，则 u 是连续的（参见 IV, 第 50 页, exerc. 15）。
\end{enumerate}

¶ 7) a) 设 E 是一个 (DF) 空间，U 是 E 中一个凸、均衡且闭的集，它吸收 E 的有界子集，并设 \((x_n)\) 是 \(\complement\) U 的点的一个序列。证明：存在 E 中 0 的一个邻域 V，它不包含任何一个 \(x_n\) 。（设 \((A_n)\) 是 E 的典范有界结构的一个可数基。证明：可以用归纳法定义一个序列 \((\lambda_n)\) （由数 \textgreater{} 0 组成）以及一个序列 \((V_n)\) （由凸、均衡且闭的

0 的邻域组成），使得 \(\lambda_{n}\mathbf{A}_{n}\subset \bigcap_{j = 1}\mathbf{V}_{j},\lambda_{n}\mathbf{A}_{n}\subset \mathrm{U}\) 且 \(x_{n}\in \mathbb{C}\mathbf{V}_{n}\) 对所有 \(n\) 。为此，若诸 \(\lambda_{j}\) 与 \(\mathrm{V}_j\) 已对

\(j \leqslant n\) 构造出来，则取 \(\lambda_{n+1}\) 使得 \(\lambda_{n+1}\mathrm{A}_{n+1} \subset \mathrm{U}\) 且 \(\lambda_{n+1}\mathrm{A}_{n+1} \subset \mathrm{V}_j\) 对所有 \(j \leqslant n\) ，然后取 \(\mathrm{V}_{n+1}\) 包含诸 \(\lambda_j\mathrm{A}_j\) （\(j \leqslant n + 1\) ）的并的凸包的闭包。

\begin{enumerate}
\def\labelenumi{\alph{enumi})}
\setcounter{enumi}{1}
\item
  由 a) 推出：若 M 是 E 的一个包含处处稠密可数集的子集，则 E 的双对偶 \(E''\) 的强拓扑在 M 上诱导的拓扑与 E 的拓扑在 M 上诱导的拓扑相同。特别地，E 中的收敛序列对于 E 的拓扑与对于 \(E''\) 的强拓扑诱导的拓扑是同样的；对 E 的每个可度量子集 M，E 的拓扑在 M 上诱导的拓扑与 \(E''\) 的强拓扑诱导的拓扑相同。
\item
  由 \(a\) ) 推出：若在 \(E\) 中存在一个可数的处处稠密集，则 \(E\) 是次桶型的。
\item
  由 b) 以及 exerc. 6 推出：若 E 的每个有界子集对于 E 的拓扑诱导的拓扑都是可度量的，则 E 是次桶型的。
\end{enumerate}

\begin{enumerate}
\def\labelenumi{\arabic{enumi})}
\setcounter{enumi}{7}
\tightlist
\item
  设 E 是一个 Fréchet 空间，\(E_{b}^{\prime}\) 是它的强对偶。假设存在一个处处稠密序列 \((x_{n}^{\prime})\) 于 \(E_{b}^{\prime}\) 中。证明：E 满足第一可数公理。（设 \((\mathbf{K}_{n}^{\prime})\) 是 \(E_{b}^{\prime}\) 的典范有界结构的一个由闭凸均衡集组成的可数基。对每个系统 \(\alpha\) ——它由一点 \(x_{n}^{\prime}\) 、任意有限个有理数 \(\lambda_{k} > 0 (1 \leqslant k \leqslant m)\) 以及 m 个指标 \(n_{k}\) 组成，且满足 \(x_{n}^{\prime} \notin 2 \sum_{k=1}^{m} \lambda_{k} K_{n_{k}}^{\prime} = 2 H_{\alpha}^{\prime}\) ——取 \(x_{\alpha} \in E\) 使得
\end{enumerate}

方程为 \(\langle x_{\alpha}, y' \rangle = 1\) 的超平面把两个弱紧集 \(\mathrm{H}_{\alpha}'\) 与 \(x_{n}' + \mathrm{H}_{\alpha}'\) 严格分离。证明：对每个 \(x' \neq 0\) （属于 \(\mathrm{E}'\) ），存在一个系统 \(\alpha\) 使得 \(\langle x_{\alpha}, x' \rangle \neq 0\) 。为此，考虑 0 的一个邻域 \(\mathrm{V}'\) （\(\mathrm{E}_b'\) 中的），使得 \(\mathrm{V}' \cap (x' + \mathrm{V}') = \emptyset\) ，然后对每个整数 \(m\) ，取有理数 \(\lambda_m > 0\) 使得 \(\lambda_m \mathrm{K}_m' \subset \mathrm{V}'\) ；利用如下事实：并 \(\mathrm{U}' \subset \mathrm{V}'\) （由诸 \(\lambda_m \mathrm{K}_m'\) 取并）是 0 的一个邻域（exerc. 7，以及 IV, 第 58 页, exerc. 3, b)），并且存在 \(n\) 使得 \(x_n' \in x' + \mathrm{U}'\) 。

¶ 9) a) 设 E 是一个 Hausdorff 半桶型空间，M 是 E 的一个闭向量子空间，E' 是 E 的对偶。证明：E/M 是半桶型的，并且强拓扑 β(M°, E/M) 与强拓扑 β(E', E) 在 M° 上诱导的拓扑相同。（注意，只需证明 M° 中一个对于 β(E', E) 收敛于 0 的序列 (x') 对于 β(M°, E/M) 是有界的。）由此推出：若 E 是 (DF) 空间，则 E/M 也是（参见 IV, 第 63 页, exerc. 8）。

\begin{enumerate}
\def\labelenumi{\alph{enumi})}
\setcounter{enumi}{1}
\item
  设 E 是一个 Hausdorff 局部凸空间，M 是 E 的一个（未必闭的）向量子空间。证明：若 M 是半桶型空间，则强拓扑 \(\beta(\mathrm{E}^{\prime}/\mathrm{M}^{\circ}, \mathrm{M})\) 与按 \(M^{\circ}\) 从强拓扑 \(\beta(\mathrm{E}^{\prime}, \mathrm{E})\) 得到的商拓扑相同。（按 a) 的方式论证。）
\item
  证明：一个 Hausdorff 且拟完备的半桶型空间 E 是完备的（对 E 与 \(\hat{E}\) 应用 b)）。特别地，一个半桶型、半自反的空间是完备的（参见 IV, 第 52 页, exerc. 6）。
\item
  证明：半桶型（相应地，(DF)）Hausdorff 空间的完备化是半桶型的（相应地，是 (DF) 空间）。
\item
  设 \((\mathrm{E}_n)\) 是半桶型（相应地，(DF)）空间的一个序列，E 是一个向量空间，且对每个 \(n\) ，设 \(f_n\) 是从 \(\mathrm{E}_n\) 到 E 中的一个线性映射。假设 \(\dot{\mathrm{E}}\) 是诸 \(f_n(\mathrm{E}_n)\) 的并；证明：E 对于使诸 \(f_n\) 都连续的最细局部凸拓扑是半桶型的（相应地，是 (DF) 空间）（先考察 E 是诸 \(\mathrm{E}_n\) 的拓扑直和的情形）。若诸 \(\mathrm{E}_n\) 都是半自反的（相应地，自反的）且 E 是 Hausdorff 的，则 E 是半自反的（相应地，自反的）。
\end{enumerate}

\begin{enumerate}
\def\labelenumi{\arabic{enumi})}
\setcounter{enumi}{9}
\tightlist
\item
  设 E 是一个满足第一可数公理的 Fréchet 空间。证明：若在 E 的对偶 \(E'\) 中，每个对于弱拓扑 \(\sigma(E', E)\) 收敛的序列也对于强拓扑 \(\beta(E', E)\) 收敛，则 E 是一个 Montel 空间。（证明：\(E'\) 的每个有界子集对于强拓扑都是相对紧的；利用 GT, II, § 4, exerc. 6；然后利用 IV, 第 53 页, exerc. 11, b)。）
\end{enumerate}

¶ 11) 设 \((c_{mn})\) 是由 \textgreater{} 0 的数组成的一个二重序列，满足 \(c_{m,n} \leqslant c_{m+1,n}\) ，并设 E 是由所有实数序列 \(x = (x_n)\) 组成的空间，这些序列使得 \(p_m(x) = \sum_n c_{mn} |x_n| < +\infty\) ，对

每个整数 \(m\) 成立。我们给 \(E\) 赋予由诸半范数 \(p_m\) 定义的拓扑，则对于该拓扑 \(E\) 是一个满足第一可数公理的 Fréchet 空间；对偶 \(E'\) （\(E\) 的对偶）可以与由所有满足如下条件的序列 \(x' = (x_n')\) 组成的空间等同：\(\sup_{n} c_{mn}^{-1} |x_n'| < +\infty\) 至少对

一个 \(m\) 成立，且典范双线性型 \(\langle x, x' \rangle\) 等同于 \(\sum_{n} x_n x_n'\) (IV, 第 47 页,

exerc. 1, c))。假设不存在任何子序列 \((n_k)\) ，使得有一个序列 \((a_m)\) （由 \(\geqslant 0\) 的数组成）以及一个指标 \(m_0\) ，满足 \(c_{m,n_k} \leqslant a_m c_{m_0,n_k}\) 对所有 \(m \geqslant m_0\) 以及所有 \(k\) 。在这些条件下，\(\mathbf{E}'\) 中每个弱收敛序列都强收敛，从而 (IV, 第 60 页, exerc. 10) E 是一个 Montel 空间。（用反证法；若有必要，作一个形如 \((x_n) \mapsto (a_n x_n)\) 的变换，化归为 \(c_{m_0 n} = 1\) 对所有 \(n\) 及某个 \(m_0\) 成立的情形，且存在一个序列 \((x'^{(p)})_{p \geqslant 0}\) 在 \(\mathbf{E}'\) 中弱收敛于 0，并满足 \(|x_n'^{(p)}| \leqslant 1\) 对每一对 \((p, n)\) ，同时存在 E 中一个有界集 B，由 \(p_m(x) \leqslant b_m\) （对所有 \(m \geqslant 0\) ）定义，且使得 \(\sup_{x \in \mathbf{B}} |\langle x, x'^{(p)} \rangle| \geqslant 2\delta > 0\) 对

每个整数 \(p\) 成立。在这些假设下，证明存在整数的一个严格递增序列 \((r_q)\) ，以及点的一个序列 \((x^{(q)})\) （\(\mathbf{B}\) 的点），使得

\[
\sum_ {k = r _ {q} + 1} ^ {r _ {q+1}} \left| x _ {k} ^ {(q)} \right| > \delta
\]

对每个指标 \(q\) 成立。然后（用反证法）证明：对每个 \(q\) ，至少存在一个指标 \(s_q\) 使得 \(r_q < s_q \leqslant r_{q+1}\) ，并且对每个整数 \(m\) 有 \(c_{m,s_q} \leqslant b_m 2^{m+2}/\delta\) ，这与假设矛盾。）

\begin{enumerate}
\def\labelenumi{\arabic{enumi})}
\setcounter{enumi}{11}
\tightlist
\item
  \begin{enumerate}
  \def\labelenumii{\alph{enumii})}
  \tightlist
  \item
    设 F 是一个 Hausdorff (DF) 空间，\(F_{b}^{\prime}\) 是它的强对偶。证明：若 \(F_{b}^{\prime}\) 是自反的，则 F 的完备化 \(\hat{F}\) 是自反的且等于 F 的双对偶 \(F^{\prime\prime}\) （参见 IV, 第 52 页, exerc. 4 及第 53 页, exerc. 11）。
  \end{enumerate}
\end{enumerate}

\begin{enumerate}
\def\labelenumi{\alph{enumi})}
\setcounter{enumi}{1}
\tightlist
\item
  设 E 是一个 Fréchet 空间。证明：若 E 的双对偶 \(E''\) 是自反的，则 E 是自反的。
\end{enumerate}

\begin{enumerate}
\def\labelenumi{\arabic{enumi})}
\setcounter{enumi}{12}
\tightlist
\item
  \begin{enumerate}
  \def\labelenumii{\alph{enumii})}
  \tightlist
  \item
    设 E、F 是两个 Fréchet 空间，G 是一个 Hausdorff 局部凸空间，\(E'\) ，\(F'\) ，\(G'\) 分别是 E、F、G 的对偶。设 u 是从 \(E' \times F'\) 到 \(G'\) 中的一个双线性映射，当 \(E'\) ，\(F'\) ，\(G'\) 分别带上弱拓扑 \(\sigma(E', E)\) ，\(\sigma(F', F)\) 与 \(\sigma(G', G)\) 时，u 是分别连续的 (III, 第 28 页) 。证明：在这些条件下，u 是从 \(E' \times F'\) 到 \(G'\) 中的一个连续映射（当 \(E'\) ，\(F'\) 与 \(G'\) 带上强拓扑时）。（对 \(z \in G\) ，令 \(\langle z, u(x', y') \rangle = \langle v_z(x'), y' \rangle\) ，其中 \(v_z(x') \in F\) 。首先证明：若 \(E'\) 带上强拓扑而 \(F\) 带上初始拓扑，则诸 \(v_z\) （当 \(z\) 取遍一个有界集 \(C\) （\(G\) 中的）时）所成的集是等度连续的；为此利用 IV, 第 51 页, exerc. 19, d)。其次证明：存在 0 的一个邻域 \(V'\) （对于 \(E'\) 的强拓扑），使得诸集 \(v_z(V')\) （当 \(z\) 取遍 \(C\) 时）的并在 \(F\) 中是有界的；为此利用 III, 第 47 页, exerc. 5。）
  \end{enumerate}
\end{enumerate}

\begin{enumerate}
\def\labelenumi{\alph{enumi})}
\setcounter{enumi}{1}
\tightlist
\item
  给出一个例子，以说明若假设 E 是 Fréchet 空间而 F 是 Fréchet 空间的一个严格归纳极限 (III, 第 47 页, exerc. 3)，则 a) 的结论不成立。
\end{enumerate}

\begin{enumerate}
\def\labelenumi{\arabic{enumi})}
\setcounter{enumi}{13}
\tightlist
\item
  \begin{enumerate}
  \def\labelenumii{\alph{enumii})}
  \tightlist
  \item
    设 E 是一个 Fréchet 空间，\(E'\) 是它的对偶。证明：\(E'\) 带上紧收敛拓扑或任何更细的 S-拓扑都是完备的（参见 III, 第 22 页, 注记 1）。若 E 不是自反的，证明：\(E'\) 对于任何比紧收敛拓扑更细且比 \(\tau(\mathrm{E}', \mathrm{E})\) 更粗的 S-拓扑都不是次桶型的。
  \end{enumerate}
\end{enumerate}

\begin{enumerate}
\def\labelenumi{\alph{enumi})}
\setcounter{enumi}{1}
\tightlist
\item
  设 \((\mathrm{E}_{\alpha})_{\alpha\in\mathbf{A}}\) 是一族 Fréchet 空间，E 是一个向量空间，且对每个 \(\alpha\in A\) ，设 \(h_{\alpha}\) 是从 \(E_{\alpha}\) 到 E 中的一个线性映射。假设 E 带上使诸 \(h_{\alpha}\) 都连续的最细局部凸拓扑 (II, 第 27 页) 后是 Hausdorff 的。证明：E 的对偶 \(E'\) 带上紧收敛拓扑或任何更细的 S-拓扑都是完备的（参见 III, 第 20 页, 定理 1）。
\end{enumerate}

\begin{enumerate}
\def\labelenumi{\arabic{enumi})}
\setcounter{enumi}{14}
\tightlist
\item
  设 \(\mathrm{E}\) 是一个无限维 Banach 空间，\((a_{n})_{n\geqslant 1}\) 是 \(\mathrm{E}\) 的点的一个可数、完全且自由的族，并设 \(\mathrm{F}_n\) 是一个 \(n\) 维子空间（\(\mathrm{E}\) 中的），由诸 \(a_{m}\) （\(m\leqslant n\) ）生成。设 \(\mathrm{S}_n\) 是方程为 \(\| x\| = n\) 的球面（\(\mathrm{E}\) 中的）；在 \(\mathrm{S}_n\cap \mathrm{F}_n\) 中设 \(\mathrm{A}_n\) 是一个有限集，使得
\end{enumerate}

\(S_{n} \cap F_{n}\) 的每一点都有距它 \(\leqslant 1/n\) 的 \(A_{n}\) 中的点。证明：\(A = \bigcap_{n=1}^{\infty} A_{n}\) 与每个闭有界集的交

都是闭的，但 0 是 A 对于弱化拓扑的一个极限点。

\begin{enumerate}
\def\labelenumi{\arabic{enumi})}
\setcounter{enumi}{15}
\tightlist
\item
  设 E 是局部凸可度量空间的一个序列 \((\mathrm{E}_{p})\) 的归纳极限空间，典范映射 \(E_{p} \rightarrow E\) 都是单射，且 E 是诸 \(E_{p}\) 的像的并。证明：E 的强对偶 \(E_{b}^{\prime}\) 是可穷竭的 (III, 第 49 页, exerc. 1)。（设 \((\mathrm{U}_{j}^{p})\) 是 \(E_{p}\) 中 0 的闭、凸且均衡邻域的一个递减基本系；考虑极集 \((\mathrm{U}_{j}^{p})^{\circ}\) （\(E^{\prime}\) 中的）的有限交；利用如下事实：对每个递增序列
\end{enumerate}

\((m_{j})_{j\geqslant 1}\) ，交 \(\bigcap_{j = 1}^{\infty}(\mathrm{U}_{m_j}^j)^{\circ}\) 是 E 中 0 的某个邻域的极集，并且 \(\mathrm{E}_b^\prime\) 是完备的。）

¶ 17) a) 设 E 是一个 Hausdorff 局部凸空间，使得由 E 的相对紧集组成的有界结构具有一个可数基 (A \(_{n}\) ) (III, 第 1 页)。证明：(A \(_{n}\) ) 也是典范有界结构的一个基 (III, 第 3 页, 定义 5)。（设 C \(_{n}\) 是 n 个都等于 A \(_{n}\) 的集之和所成的相对紧集；则 (C \(_{n}\) ) 也是相对紧集有界结构的一个基。用反证法，考虑一个不包含于任何 C \(_{n}\) 中的有界集 B，推断存在 B 的点的一个序列 (x \(_{n}\) ) 使得 x \(_{n}\) /n∉ A \(_{n}\) ；由此导出矛盾。）于是，空间 E 是半自反的，并且 E 中每个紧集的闭凸均衡包都是紧的。

\begin{enumerate}
\def\labelenumi{\alph{enumi})}
\setcounter{enumi}{1}
\tightlist
\item
  假设 E 是次桶型的，并且对一个与 E 和 \(E'\) 之间的对偶相容的拓扑 T，E 中对于 T 的相对紧子集组成的有界结构具有可数基。证明：此时 E 是一个自反的 (DF) 空间 (IV, 第 57 页, exerc. 2)。（利用 a)，注意 E 中的闭有界集对于 T 都是紧的，从而对于 E 的初始拓扑是完备的；这就蕴含 E 是桶型的。）若 T 是初始拓扑，则 E 是一个 Montel 空间。
\end{enumerate}

\begin{enumerate}
\def\labelenumi{\arabic{enumi})}
\setcounter{enumi}{17}
\tightlist
\item
  \begin{enumerate}
  \def\labelenumii{\alph{enumii})}
  \tightlist
  \item
    设 E 是一个 Hausdorff 局部凸空间，\((\mathbf{A}_{n})\) 是对于弱化拓扑 \(\sigma(\mathrm{E}, \mathrm{E}')\) 凸、均衡、紧的集的一个递增序列，满足诸集 \(A_{n}\) 的并是 E，且对每个整数 n 及每个 \(\lambda > 0\) ，存在 m 使得 \(\lambda A_{n} \subset A_{m}\) 。证明：\((\mathbf{A}_{n})\) 是由对于 \(\sigma(\mathrm{E}, \mathrm{E}')\) 凸且相对紧的集组成的有界结构的一个基。（注意：在 \(E'\) 上，在诸 \(A_{n}\) 上一致收敛的拓扑就是 \(\tau(\mathrm{E}', \mathrm{E})\) 。）b) 若 E 是桶型的，且存在一个具有 a) 中所述性质的序列 \((\mathbf{A}_{n})\) ，则 E 是一个自反的 (DF) 空间。（注意：\(E'\) 带上 \(\tau(\mathrm{E}', \mathrm{E})\) 是可度量的，并且 E 的拓扑是 \(\beta(\mathrm{E}, \mathrm{E}')\) 。）
  \end{enumerate}
\end{enumerate}

\exercisegroup

\begin{enumerate}
\def\labelenumi{\arabic{enumi})}
\item
  设 E 与 F 是两个 Hausdorff 局部凸空间，\(E'\) 与 \(F'\) 分别是它们的对偶。证明：若对 F 的每个向量子空间 N，\(\tau(\mathrm{F}, \mathrm{F}')\) 在 N 上诱导的拓扑都与 \(\tau(\mathrm{N}, \mathrm{F}'/\mathrm{N}^{\circ})\) 相同 (IV, 第 51 页, exerc. 20)，则每个对于拓扑 \(\sigma(\mathrm{E}, \mathrm{E}')\) 与 \(\sigma(\mathrm{F}, \mathrm{F}')\) 的从 E 到 F 中的严格态射，也是对于拓扑 \(\tau(\mathrm{E}, \mathrm{E}')\) 与 \(\tau(\mathrm{F}, \mathrm{F}')\) 的严格态射（参见 IV, 第 10 页, 命题 11）。考察 F 可度量的情形。
\item
  设 E 是一个 Hausdorff 局部凸空间，使得在对偶 \(E'\) 中存在一个无限维凸集 \(B'\) ，它对于 \(\sigma(E', E)\) 是紧的（例如当 E 是一个带上 \(\sigma(E, E^{*})\) 的无限维向量空间时，该条件即成立）。证明：存在一个线性型 \(u \in E'^{*}\) ，它在 \(B'\) 上是无界的；由此推出 \(B'\) 对于拓扑 \(\sigma(E', F)\) 不是紧的，其中 F 是子空间 \(E + Ru\) （\(E'^{*}\) 的）。由此推断：从 E 到 F 的典范单射对于拓扑 \(\sigma(E, E')\) 与 \(\sigma(F, E')\) 是一个严格态射，但对于拓扑 \(\tau(E, E')\) 与 \(\tau(F, E')\) 不是。
\item
  给出一个严格单射态射 \(u\) （从 Fréchet 空间 \(E\) 到 Fréchet 空间 \(F\) 中）的例子，使得 \(^t u\) 不是从 \(F_b'\) 到 \(E_b'\) 中的严格态射（参见 IV, 第 58 页, exerc. 5, c)）。
\item
  设 E 与 F 是两个 Hausdorff 局部凸空间，u 是从 E 到 F 中的一个连续线性映射，M 是 E 的一个处处稠密的向量子空间。证明：若 u 在 M 上的限制是从 M 到 F 中的一个严格态射，则 u 是从 E 到 F 中的一个严格态射（利用 IV, 第 27 页的命题 2）。若此外还有 \(u(\mathbf{M}) = \mathbf{F}\) ，证明：对 E 中 0 的每个开的、凸且均衡的邻域 V，\(u(\mathbf{V})\) 都属于 \(\overline{u(\mathbf{V} \cap \mathbf{M})}\) 的内部。
\item
  设 \(\mathbf{E}\) 与 \(\mathbf{F}\) 是两个赋范空间，\(u\) 是从 \(\mathbf{E}\) 到 \(\mathbf{F}\) 中的一个连续线性映射。
\end{enumerate}

\begin{enumerate}
\def\labelenumi{\alph{enumi})}
\item
  证明：若 \(u\) 是从 \(E\) 到 \(F\) 中的一个严格态射，则 \({}^t u\) 是从强对偶 \(F_b'\) 到强对偶 \(E_b'\) 中的一个严格态射。
\item
  假设 E 是完备的；证明：若 \(^{t}u\) 是从 \(F_{b}^{\prime}\) 到 \(E_{b}^{\prime}\) 中的一个严格态射，则 u 是从 E 到 F 中的一个严格态射，且 \(^{t}u\) 是从 \(F^{\prime}\) 到 \(E^{\prime}\) 中的一个严格态射（对于弱拓扑 \(\sigma(F^{\prime}, F)\) 与 \(\sigma(E^{\prime}, E)\) ）（把 F 视为它的完备化的一个子空间）。
\item
  给出一个例子：E 不是完备的，F 是完备的，\(^{t}u\) 对于强拓扑以及对于弱拓扑都是从 \(F'\) 到 \(E'\) 中的一个严格单射态射，但 u 不是从 E 到 F 中的严格态射（参见 II, 第 74 页, exerc. 5）。
\item
  给出一个例子：E 不是完备的，F 是完备的，u 是从 E 到 F 中的一个严格单射态射，\(^{t}u\) 对于强拓扑是从 F' 到 E' 中的一个严格态射，但对于弱拓扑不是（取 E 在 F 中处处稠密）。
\item
  若 F 是完备的，且 \({}^t u\) 对于弱拓扑是从 \(\mathrm{F}'\) 到 \(\mathrm{E}'\) 中的一个严格态射，则 \({}^t u\) 对于强拓扑是从 \(\mathrm{F}'\) 到 \(\mathrm{E}'\) 中的一个严格态射（把 \(u\) 延拓到 \(\hat{\mathrm{E}}\) 上）。
\item
  给出一个例子：E 是完备的，F 不是完备的，\(^{t}u\) 对于弱拓扑是从 \(F'\) 到 \(E'\) 中的一个严格态射，但对于强拓扑不是（参见 II, 第 74 页, exerc. 5）。
\end{enumerate}

\begin{enumerate}
\def\labelenumi{\arabic{enumi})}
\setcounter{enumi}{5}
\item
  设 E 是可度量的局部凸空间 \(\ell^{1}(\mathbf{N})\) (I, 第 4 页) ，带上由乘积空间 \(R^{N}\) 的拓扑诱导的拓扑；它的对偶 \(E'\) 可以与 \(\mathbf{R}^{(N)}\) 等同，并且拓扑 \(\tau(\mathrm{E}',\mathrm{E})\) 是在 \(E'\) 上由 \(c_{0}(\mathbf{N})\) 的范数拓扑诱导的拓扑 (IV, 第 47 页, exerc. 1)。证明：若 u 是从 E 到一个 Hausdorff 桶型空间 F 上的满射连续线性映射，则 F 必是有限维的。（注意：\({}^{t}u\) 是从 \(F'\) （带上 \(\sigma(\mathrm{F}',\mathrm{F})\) ）到 \(E'\) （带上 \(\sigma(\mathrm{E}',\mathrm{E})\) ）的一个子空间上的同构；由 F 是桶型的这一事实，推出 \({}^{t}u(F')\) 带上由 \(\tau(\mathrm{E}',\mathrm{E})\) 诱导的拓扑是一个 Banach 空间，然后利用 II, 第 80 页的 exerc. 24。）
\item
  \begin{enumerate}
  \def\labelenumii{\alph{enumii})}
  \tightlist
  \item
    设 \(\mathbf{E}\) 是一个 Banach 空间，\((x_{\alpha})_{\alpha \in \mathbf{A}}\) 是 \(\mathbf{E}\) 的单位球面中的一个处处稠密集。设 \(u\) 是从空间 \(\ell^1 (\mathbf{A})(\mathrm{I},\mathrm{p.~4})\) 到 \(\mathbf{E}\) 中的线性映射，由 \(u(t) = \sum_{\alpha \in \mathbf{A}}t(\alpha)x_{\alpha}\) 定义。
  \end{enumerate}
\end{enumerate}

对所有 \(t = (t(\alpha))_{\alpha \in \mathbf{A}}\) （属于 \(\ell^1 (\mathbf{A})\) ）成立。证明：\(u\) 是从 \(\ell^1 (\mathbf{A})\) 到 \(\mathbf{E}\) 上的一个严格态射，从而 \(\mathbf{E}\) 同构于 \(\ell^1 (\mathbf{A})\) 的一个商空间。

\begin{enumerate}
\def\labelenumi{\alph{enumi})}
\setcounter{enumi}{1}
\tightlist
\item
  由 a) 出发，给出 \(\ell^{1}(\mathbf{N})\) 的一个在 \(\ell^{1}(\mathbf{N})\) 中没有拓扑补的闭子空间的例子（对 E 取 \(c_{0}(\mathbf{N})\) ，并利用 IV, 第 54 页, exerc. 15, b) 与第 55 页, exerc. 18）。
\end{enumerate}

\begin{enumerate}
\def\labelenumi{\arabic{enumi})}
\setcounter{enumi}{7}
\tightlist
\item
  对每个整数 \(n > 0\) ，设 \(a^{(n)}\) 是二重序列 \(a^{(n)} = (a_{ij}^{(n)})_{i\geqslant 1,j\geqslant 1}\) ，其定义为：\(a_{ij}^{(n)} = j^n\) 对每一对 \((i,j)\) （满足 \(i < n\) ），而 \(a_{ij}^{(n)} = i^n\) 对每一对 \((i,j)\) （满足 \(i\geqslant n\) ）。设 \(E\) 是由所有满足如下条件的实数二重序列 \(x = (x_{ij})\) 组成的向量空间：对每个整数 \(n > 0\) ，有 \(p_n(x) = \sum_{i,j}a_{ij}^{(n)}|x_{ij}| < +\infty\) ；诸半范数 \(p_n\) 在 E 上定义
\end{enumerate}

一个 Fréchet 空间兼 Montel 空间的拓扑 (IV, 第 60 页, exerc. 11)；E 的对偶 E' 是一个 (DF) 空间兼 Montel 空间（从而是超有界型且自反的），它等同于由所有满足如下条件的序列 \(x' = (x'_{ij})\) 组成的空间：至少对一个指标 n，有关系 \(\sup |a_{ij}^{(n)}|^{-1} |x'_{ij}| < +\infty\) (IV, 第 47 页, exerc. 1, c))。

\begin{enumerate}
\def\labelenumi{\alph{enumi})}
\item
  对每个 \(x = (x_{ij}) \in E\) ，令 \(y_j = \sum_i x_{ij}\) （对所有 \(j \geqslant 1\) ）。证明：\(\sum_j |y_j| < +\infty\) ；设 \(u(x)\) 表示序列 \((y_i) \in \ell^1(\mathbf{N})\) ；证明 u 是从 E 到 \(F = \ell^1(\mathbf{N})\) 中的一个连续线性映射，并且对每个序列 \(y' = (y'_j) \in F' = \ell^\infty(\mathbf{N})\) ，\(^t u(y')\) 是序列 \((z'_{ij}) \in E'\) ，它由 \(z'_{ij} = y'_j\) （对每个指标 i）确定。由此推出：\(^t u\) 是从 \(\ell^\infty(\mathbf{N})\) 到 \(E'\) 的一个对于 \(\sigma(E', E)\) 闭的子空间上的单射线性映射，从而 u 对于初始拓扑是从 E 到 \(\ell^1(\mathbf{N})\) 上的一个严格态射，且 \(^t u\) 是从 \(F' = \ell^\infty(\mathbf{N})\) 到 \(^t u(F')\) 上的一个同构（对于弱拓扑 \(\sigma(F', F)\) 与 \(\sigma(E', E)\) ）。
\item
  设 \(M = u^{-1}(0)\) ；则 \(M^{\circ} = {}^{t}u(F')\) 。证明：在 \({}^{t}u\) 下，\(M^{\circ}\) 上由 \(E'\) 的强拓扑诱导的拓扑的原像，是 F 的紧子集上一致收敛的拓扑 (IV, 第 28 页, 定理 1 及第 54 页, exerc. 15)。由此推出：在 \(M^{\circ}\) 上，由强拓扑 \(\beta(E', E)\) 诱导的拓扑不是强拓扑 \(\beta(M^{\circ}, E/M)\) ，并且对于由 \(\beta(E', E)\) 诱导的拓扑，\(M^{\circ}\) 不是次桶型空间，尽管它是一个超有界型 Montel 空间的闭子空间；另一方面，E/M（它是 Fréchet 空间兼 Montel 空间关于一个闭子空间的商）不是自反的。证明：E/M 中存在一些有界集，它们不是 E 的有界子集的典范像。
\end{enumerate}

\begin{enumerate}
\def\labelenumi{\arabic{enumi})}
\setcounter{enumi}{8}
\tightlist
\item
  设 A 是一个可数集。考虑三个向量空间对 \((P, P')\) ，\((Q, Q')\) ，\((E, E')\) ，这六个空间都是 \(R^{A}\) 的向量子空间，且都包含直和子空间 \(\mathbf{R}^{(A)}\) ；此外，假设对每一点 \(x \in P\) （相应地，\(x \in Q\) ，\(x \in E\) ）以及每一点 \(x' \in P'\) （相应地，\(x' \in Q'\) ，\(x' \in E'\) ），族 \((x_{\alpha}x_{\alpha}')_{\alpha \in A}\) 是可和的，并令 \(\langle x, x' \rangle = \sum_{\alpha \in A} x_{\alpha}x_{\alpha}'\) ；
\end{enumerate}

这个双线性型使 P 与 P' （相应地，Q 与 Q' ，E 与 E' ）处于分离对偶之中。

\begin{enumerate}
\def\labelenumi{\alph{enumi})}
\item
  假设 \(\mathrm{E} = \mathrm{P} \cap \mathrm{Q}\) ，\(\mathrm{E}' \supset \mathrm{P}' + \mathrm{Q}'\) 且 \(\mathrm{E}' \neq \mathrm{P}' + \mathrm{Q}'\) 。考虑线性映射 \(u: x \mapsto (x, x)\) （从 \(\mathrm{E}\) 到 \(\mathrm{F} = \mathrm{P} \times \mathrm{Q}\) 中），它与 \(\mathrm{F}' = \mathrm{P}' \times \mathrm{Q}'\) 处于分离对偶之中。证明：\(u\) 对于弱拓扑 \(\sigma(\mathrm{E}, \mathrm{E}')\) 与 \(\sigma(\mathrm{F}, \mathrm{F}')\) 是连续的，并且它的像 \(\mathbf{M} = u(\mathbf{E})\) 对于 \(\sigma(\mathbf{F}, \mathbf{F}')\) 是一个闭子空间；由此推出：\(^t u\) 是从 \(\mathbf{F}'\) 到 \(\mathbf{E}'\) 中的一个严格态射（对于拓扑 \(\sigma(\mathbf{F}', \mathbf{F})\) 与 \(\sigma(\mathbf{E}', \mathbf{E})\) ），而 \(\mathbf{N} = {}^{t}u(\mathbf{F}')\) 在 \(\mathbf{E}'\) 中对于 \(\sigma(\mathbf{E}', \mathbf{E})\) 不是闭的。若 \(\mathbf{E}'\) 对于拓扑 \(\tau(\mathbf{E}', \mathbf{E})\) 是可度量的，推出：\(^t u\) 也是从 \(\mathbf{F}'\) 到 \(\mathbf{E}'\) 中的一个严格态射（对于拓扑 \(\tau(\mathbf{F}', \mathbf{F})\) 与 \(\tau(\mathbf{E}', \mathbf{E})\) ）。
\item
  此外，假设 \(E'\) 对于 \(\tau(E', E)\) 是一个 Fréchet 空间。则 \(F'/M^{\circ}\) 带上 \(\tau(F', F)\) 关于 \(M^{\circ}\) 的商拓扑后不是半完备的，并且 \(F'/M^{\circ}\) 中存在一些对于 \(\tau(F'/M^{\circ}, M)\) 不是相对紧的有界集。
\item
  在与 b) 相同的假设下，设 \(x'\) 是 \(E'\) 的一个不属于 \(\mathbf{N} = {}^{u}(\mathbf{F}')\) 的元素；对每个 \(y \in M\) ，令 \(v(y) = \langle x, x' \rangle\) ，其中 \(x \in E\) 是使 \(u(x) = y\) 的唯一元素。证明：M 上的线性型 v 对于拓扑 \(\sigma(F, F')\) 不是连续的，但它在 M 的每个有界子集上的限制对于 \(\sigma(M, F'/M^{\circ})\) 是连续的。由此推出：\(L = v^{-1}(0)\) 是 F 的一个向量子空间，它与 F 的每个有界闭子集（对于 \(\sigma(F, F')\) ）的交对于 \(\sigma(F, F')\) 是闭的，但它本身在 F 中对于 \(\sigma(F, F')\) 不是闭的。
\end{enumerate}

¶ 10) a) 设 G、H 是两个自反 Banach 空间，满足 \(\mathbf{R}^{(\mathbf{N})} \subset \mathbf{G} \subset \mathbf{H} \subset \mathbf{R}^{\mathbf{N}}\) （* 例如 \(\mathbf{G} = \ell^{r}(\mathbf{N})\) 且 \(\mathbf{H} = \ell^{p}(\mathbf{N})\) ，其中 \(1 < r < p < +\infty\) ）。取 \(\mathbf{A} = \mathbf{N} \times \mathbf{N}\) ，\(\mathbf{P} = \mathbf{H}^{(\mathbf{N})}\) （拓扑直和），\(\mathbf{P}' = \mathbf{H}'^{\mathbf{N}}\) ，\(\mathbf{Q} = \mathbf{G}^{\mathbf{N}}\) ，\(\mathbf{Q}' = \mathbf{G}'^{(\mathbf{N})}\) ，\(\mathbf{E} = \mathbf{G}^{(\mathbf{N})}\) ，\(\mathbf{E}' = \mathbf{G}'^{\mathbf{N}}\) ，证明：exerc. 9, a) 的条件满足；\(\mathbf{E}'\) 是一个自反 Fréchet 空间，且 E、F、F' 都是自反 Fréchet 空间的严格归纳极限（从而是完备且自反的）。

\begin{enumerate}
\def\labelenumi{\alph{enumi})}
\setcounter{enumi}{1}
\tightlist
\item
  由 a) 与 exerc. 9 出发，构造具有下列性质的例子：
\end{enumerate}

\(\alpha)\) 自反 Fréchet 空间的严格归纳极限的一个商空间，它既非拟完备亦非半自反。

β) 自反 Fréchet 空间的严格归纳极限的一个闭子空间，它不是自反的，且其对偶不是完备的。

\(\gamma)\) 自反 Fréchet 空间的严格归纳极限 \(\mathbf{E}_n\) 的一个向量子空间，它不是闭的（从而不是桶型的），但它与每个子空间 \(\mathbf{E}_n\) 的交都是闭的。

δ) 自反 Fréchet 空间的严格归纳极限的对偶的一个非闭子空间，它与每个弱紧子集的交都是弱紧的（参见 IV, 第 25 页, 推论 2）。

\[
\mathrm{E} _ {n}
\]

\[
\mathrm{F} _ {n} (c f.
\]

\[
\mathrm{F} _ {n} = \overline {{\mathrm{F} \cap \mathrm{E} _ {n}}}
\]

\[
u _ {n}
\]

\[
v _ {n}
\]

\[
\mathrm{F} \cap \mathrm{E} _ {n}
\]

\[
\overline {{\mathrm{F} \cap \mathrm{E} _ {n}}}
\]

\[
v _ {n}
\]

\begin{enumerate}
\def\labelenumi{\alph{enumi})}
\setcounter{enumi}{1}
\tightlist
\item
  假设 F 在 E 中不是闭的，但对所有 n，\(F \cap E_{n}\) 在 \(E_{n}\) 中是闭的 (exerc. 10, b))。设 \(A_{n}\) 是 \(E_{n}\) 中的一个可数稠密集，G 是 E 中由 F 与诸 \(A_{n}\) 的并所生成的向量子空间。证明：G 是一个有界型空间，其中 F 是一个具有可数基的补的子空间，但 F 不是次桶型的（参见 III, 第 41 页, § 2, exerc. 4 及第 45 页, exerc. 17）。
\end{enumerate}

\begin{enumerate}
\def\labelenumi{\arabic{enumi})}
\setcounter{enumi}{11}
\tightlist
\item
  设 E、F 是两个 Fréchet 空间，u 是从 E 到 F 中的一个严格态射。证明：对每个从 E 到 F 中的有限秩连续线性映射 v，\(u + v\) 都是从 E 到 F 中的一个严格态射。
\end{enumerate}

¶ 13) a) 设 E 是一个 Hausdorff 且完备的局部凸空间。假设存在 E 中闭向量子空间的一个递减序列 ( \(F_{n}\) )，使得对 E 中 0 的每个邻域 V，存在 n 使得 \(F_{n} \subset V\) 。证明：空间 E 是极小型的 (II, 第 85 页, exerc. 13)。

\begin{enumerate}
\def\labelenumi{\alph{enumi})}
\setcounter{enumi}{1}
\item
  设 E 是一个满足第一可数公理且非极小型的 Fréchet 空间。证明：E 中存在两个闭向量子空间 M、N，使得 \(M \cap N = \{0\}\) 且 \(M + N\) 不是闭的。（设 \((x_{n}')\) 是 E 上线性无关的连续线性型的一个序列，构成 \(\sigma(E', E)\) 下的一个完全集 (III, 第 19 页, 推论 2)；设 \(L_{n}\) 是 E 中与诸 \(x_{i}'\) （指标 \(i \leqslant 2n\) ）正交的子空间；设 \(x_{n}, y_{n}\) 是 \(L_{n+1}\) 关于 \(L_{n}\) 的余中的两个线性无关向量。设 d 是定义 E 的拓扑的一个平移不变距离。利用 a)，证明存在一个数 \(\alpha > 0\) ，使得可取 \(d(0, x_{n}) \geqslant \alpha\) ，\(d(0, y_{n}) \geqslant \alpha\) 且 \(d(x_{n}, y_{n}) \leqslant 1/n\) 。证明：若 M（相应地，N）是由诸 \(x_{n}\) （相应地，\(y_{n}\) ）生成的闭向量子空间，则 M 与 N 具有所要求的性质；利用 I, 第 19 页, 推论 4。）
\item
  设 \(\mathbf{E}\) 是赋范空间的一个乘积 \(\prod_{\alpha \in \mathbf{A}} \mathrm{F}_{\alpha}\) 的闭子空间；证明：若 E 中由点的可数族生成的每个闭
\end{enumerate}

子空间都是极小型的，则 E 本身是极小型的。（用反证法，假设 E 在每个 \(F_{\alpha}\) 上的投影都等于 \(F_{\alpha}\) ，且对某个 \(\alpha\) ，\(F_{\alpha}\) 是无限维的。）

\begin{enumerate}
\def\labelenumi{\alph{enumi})}
\setcounter{enumi}{3}
\tightlist
\item
  设 E 是一个 Hausdorff 且完备的局部凸空间，其中由点的可数族生成的每个子空间都是可度量的。要使 E 是极小型的，必须且只需对 E 的每对满足 \(M \cap N = \{0\}\) 的闭向量子空间 M、N，\(M + N\) 在 E 中是闭的（利用 b) 与 c)）。
\end{enumerate}

\begin{enumerate}
\def\labelenumi{\arabic{enumi})}
\setcounter{enumi}{13}
\tightlist
\item
  \begin{enumerate}
  \def\labelenumii{\alph{enumii})}
  \tightlist
  \item
    设 E 是一个不存在连续范数的 Fréchet 空间。证明：E 中存在一个极小型的闭向量子空间 (II, 第 85 页, exerc. 13)，它是无限维的，从而在 E 中具有一个拓扑补。（存在 0 的凸、均衡邻域的一个严格递减基本序列 \((\mathrm{V}_{n})\) ，以及 E 的点的一个序列 \((x_{n})\) ，使得 \(x_{n} \notin V_{n+1}\) ，且过 \(x_{n}\) 的直线包含于 \(V_{n}\) 中，诸 \(x_{n}\) 线性无关；证明所求的空间是由诸 \(x_{n}\) 生成的闭向量子空间。）
  \end{enumerate}
\end{enumerate}

\begin{enumerate}
\def\labelenumi{\alph{enumi})}
\setcounter{enumi}{1}
\item
  设 E 是一个 Fréchet 空间，其拓扑不能由单独一个范数、而由范数的一个递增序列 \((p_{n})\) 定义。设 \(V_{n}\) 是由 \(p_{n}(x) \leqslant 1\) 定义的邻域，并设 \(A_{n}^{\prime} = V_{n}^{\circ}\) 是它在 \(E^{\prime}\) 中的极集；可以假设 \(A_{n+1}^{\prime}\) 不包含于 \(E^{\prime}\) 的、由 \(A_{n}^{\prime}\) 生成的向量子空间中 (IV, 第 49 页, exerc. 12, c))。设 \((x_{n}^{\prime})\) 是 \(E^{\prime}\) 的点的一个序列，满足 \(x_{n}^{\prime} \in A_{n}^{\prime}\) 且 \(x_{n}^{\prime}\) 不属于由 \(A_{n-1}^{\prime}\) 生成的向量子空间；证明：向量子空间 \(M^{\prime}\) （由序列 \((x_{n}^{\prime})\) 生成）在 \(E^{\prime}\) 中是弱闭的（参见 IV, 第 25 页, 推论 2），并且在 \(E^{\prime}\) 中对于拓扑 \(\sigma(E^{\prime}, E)\) 没有拓扑补。（注意：若存在一个弱连续投影 \(u^{\prime}\) （从 \(E^{\prime}\) 到 \(M^{\prime}\) 上），则由 Baire 定理，\(u^{\prime}(A_{1}^{\prime})\) 将包含于诸集 \(M^{\prime} \cap A_{n}^{\prime}\) 之一中，由此导出矛盾，因为 \(A_{1}^{\prime}\) 在 \(E^{\prime}\) 中是弱完全的。）由此推出：E 的子空间 \(M^{\circ}\) 在 E 中没有拓扑补。
\item
  设 E 是一个其拓扑不能由单独一个范数定义的 Fréchet 空间。证明：若 E 不同构于 Banach 空间的一个乘积，则 E 中存在一个没有拓扑补的闭向量子空间。（用反证法；设 \((p_{n})_{n\geqslant1}\) 是 E 上定义 E 的拓扑的半范数的一个递增序列；设 \(F_{n}=p_{n}^{-1}(0)\) ，并设 \(E_{n+1}\) 是 \(F_{n+1}\) 关于 \(F_{n}\) 的一个拓扑补（令 \(E=F_{0}\) ）；利用 b)，证明 \(E_{n+1}\) 是一个 Banach 空间，且 E 同构于诸 \(E_{n}(n\geqslant1)\) 的乘积；为此利用 I, 第 17 页, 定理 1。）
\end{enumerate}

\begin{enumerate}
\def\labelenumi{\arabic{enumi})}
\setcounter{enumi}{14}
\tightlist
\item
  \begin{enumerate}
  \def\labelenumii{\alph{enumii})}
  \tightlist
  \item
    设 E 是一个无限维赋范空间（实的或复的）。证明：E 中存在一个序列 \((x_{n})\) ，使得对标量的每个有界序列 \((\lambda_{n})\) ，存在 E 上的一个连续线性型 \(x^{\prime}\) ，使得 \(\langle x_n,x'\rangle = \lambda_n\) 对所有 \(n\) 成立。（构造 E 的对偶 E' 中点的一个序列 \((x_n')\) 以及 E 中点的一个序列 \((x_{n})\) ，使得 \(\langle x_i,x_j'\rangle = \delta_{ij}\) 且 \(\| x_n'\| \leqslant 2^{-n}\) 。）
  \end{enumerate}
\end{enumerate}

\begin{enumerate}
\def\labelenumi{\alph{enumi})}
\setcounter{enumi}{1}
\tightlist
\item
  要使可度量的局部凸空间 \(\mathbf{E}\) 具有 \(a\) ) 中所述的性质，必须且只需 \(\mathbf{E}\) 的完备化不是极小型空间 (II, 第 85 页, exerc. 13)。
\end{enumerate}

\begin{enumerate}
\def\labelenumi{\arabic{enumi})}
\setcounter{enumi}{15}
\tightlist
\item
  \begin{enumerate}
  \def\labelenumii{\alph{enumii})}
  \tightlist
  \item
    设 E 与 F 是两个无限维复赋范空间，u 是从 E 到 F 上的一个双射半线性映射，它相对于 C 的一个自同构 \(\sigma\) ，并且把 E 的每个闭超平面变成 F 的闭超平面。证明：C 的自同构 \(\sigma\) 必是连续的（从而或者是恒同，或者是自同构 \(\xi \mapsto \overline{\xi}\) ）。（用反证法：设 \((x_{n})\) 是 E 中满足 exerc. 15, a) 的条件的点的一个序列，并设 \((\lambda_{n})\) 是复数的一个有界序列，使得对所有 n 有 \(|\lambda_{n}^{\sigma}| \geqslant n\) . \(\|u(x_{n})\|\) ；若 \(x' \in E'\) 使得对所有 n 有 \(\langle x_{n}, x' \rangle = \lambda_{n}\) ，考虑 u 在闭超平面（由所有 \(x \in E\) （满足 \(\langle x, x' \rangle = 1\) ）组成）上的像，并导出矛盾。）
  \end{enumerate}
\end{enumerate}

\begin{enumerate}
\def\labelenumi{\alph{enumi})}
\setcounter{enumi}{1}
\tightlist
\item
  由 a) 推出：u 是从 E 到 F 中的一个连续半线性映射 (IV, 第 7 页, 推论)。
\item
  设 \(\sigma\) 是域 C 的一个不连续自同构。证明：双射 \((\xi_{n})\mapsto(\xi_{n}^{\sigma})\) （从 \(C^{N}\) 到其自身上的）把每个闭超平面变成闭超平面（参见 IV, 第 14 页, 命题 15）。
\end{enumerate}

17) 设 E、F 是两个 Banach 空间，u 是从 E 到 F 上的一个严格态射。则存在
一个数 \(m > 0\) ，使得对每个 \(\varepsilon \in ]0, m[\) 以及每一点 \(y \in F\) ，存在 \(z \in E\) ，使
得 \(u(z) = y\) 且 \(\| z \| \leqslant (m - \varepsilon)^{-1} \| y \|\)。
a) 设 \(\mathbf{B}\) 是 E 中的一个闭球 \(\| x - a \| \leqslant r\) ，并设 w 是从 \(\mathbf{B}\) 到 F 中的一个映射，满足
下列条件：\(1^{\circ}\ \sup_{x \in \mathbf{B}} \| w(x) \| = M < +\infty; 2^{\circ}\) 存在一个数 \(k > 0\) ，
使得对所有 \(x, x'\) （属于 \(\mathbf{B}\) ），有 \(\| w(x) - w(x') \| \leqslant k \| x - x' \|\) （“Lipschitz 条件”）。证明：若 \(k < m\) 且 \(M < r(m - k)\) ，则 \(\mathbf{B}\) 在映射 \(x \mapsto u(x) + w(x)\)
下的像包含一个以 \(b = u(a)\) 为中心的球。（证明：对每个 \(y \in F\) （足够接近于 \(b\) 的），可以定义
一个序列 \((x_n)_{n \geqslant 0}\) （\(\mathbf{B}\) 的点的），使得 \(u(x_0) = y\) 且 \(u(x_n) = y - w(x_{n-1})\) （对 \(n \geqslant 1\) ），
并且使得 \((x_n)\) 收敛于 \(\mathbf{B}\) 的一点。）


\begin{enumerate}
\def\labelenumi{\alph{enumi})}
\setcounter{enumi}{1}
\tightlist
\item
  假设 \(w\) 是从整个 \(E\) 到 \(F\) 中的一个映射，使得 \(\| w(x) - w(x') \| \leqslant k \| x - x' \|\) 对所有 \(x, x'\) （属于 \(E\) ）成立。证明：若 \(k < m\) ，则 \(u + w\) 是从 \(E\) 到 \(F\) 上的一个满射且开的映射。c) 特别地，若 \(w\) 是从 \(E\) 到 \(F\) 中的一个连续线性映射且 \(\| w \| < m\) ，则 \(v = u + w\) 是从 \(E\) 到 \(F\) 上的一个严格态射。若 \(\varepsilon \in ]0, m[\) ，则对每个 \(x \in v^{-1}(0)\) （满足 \(\| x \| = 1\) ），存在 \(x_0 \in u^{-1}(0)\) 使得 \(\| x - x_0 \| \leqslant \| w \| / (m - \varepsilon)\) ；若此外 \(\| w \| < m - \varepsilon\) ，则对每个 \(x_0 \in u^{-1}(0)\) （满足 \(\| x_0 \| = 1\) ），存在 \(x \in v^{-1}(0)\) 使得
\end{enumerate}

\[
\left\| x - x _ {0} \right\| \leqslant \| w \| / (m - \varepsilon - \| w \|).
\]

由此推出：若存在 E 的一个闭向量子空间 G，它是 \(u^{-1}(0)\) 的补，则 G 也是 \(v^{-1}(0)\) 的补，只要 \(\|w\|\) 足够小（用反证法：证明 \(v^{-1}(0)\) 到 \(u^{-1}(0)\) 上的投影不能包含于 \(u^{-1}(0)\) 的一个闭超平面中；另一方面，注意存在 a > 0，使得每个点 \(x \in G\) （满足 \(\|x\| = 1\) ）都有 \(\geqslant a\) 的距离到 \(u^{-1}(0)\) ）。

\begin{enumerate}
\def\labelenumi{\arabic{enumi})}
\setcounter{enumi}{17}
\tightlist
\item
  \begin{enumerate}
  \def\labelenumii{\alph{enumii})}
  \tightlist
  \item
    设 E、F 是两个赋范空间，u 是从 E 到 F 中的一个严格单射态射；则存在一个数 m > 0，使得 \(\|u(x)\| \geqslant m \|x\|\) 对所有 \(x \in E\) 成立。证明：若 w 是从 E 到 F 中的一个连续线性映射且 \(\|w\| < m\) ，则 \(v = u + w\) 是从 E 到 F 中的一个严格单射态射。此外，对每个 \(y_0 \in u(E)\) （满足 \(\|y_0\| = 1\) ），存在 \(y \in v(E)\) 使得 \(\|y - y_0\| \leqslant \|w\| / m\) ；对每个 \(y \in v(E)\) （满足 \(\|y\| = 1\) ），存在 \(y_0 \in u(E)\) 使得 \(\|y - y_0\| \leqslant \|w\| / (m - \|w\|)\) 。
  \end{enumerate}
\end{enumerate}

\begin{enumerate}
\def\labelenumi{\alph{enumi})}
\setcounter{enumi}{1}
\tightlist
\item
  由 a) 推出：若此外 E 与 F 都是 Banach 空间，且存在 F 的一个闭向量子空间 G，它是闭子空间 \(u(\mathrm{E})\) 的补，则 G 也是 \(v(\mathrm{E})\) 的补，只要 \(\|w\|\) 足够小（按 exerc. 17, c) 的方式论证）。
\end{enumerate}

\begin{enumerate}
\def\labelenumi{\arabic{enumi})}
\setcounter{enumi}{18}
\item
  设 E 是 Banach 空间 \(\mathcal{C}(\left[-1,1\right];\mathbf{R})\) （由 \([-1,1]\) 到 R 中的全体连续映射组成）中由多项式组成的子空间。类似地，设 F 是 \(\mathcal{C}(\left[0,1\right];\mathbf{R})\) 中由全体多项式组成的子空间。设 u 是使每个多项式 \(f\in E\) 对应于 F 中的多项式 \(t\mapsto\frac{1}{2}(f(\sqrt{t})+f(-\sqrt{t}))\) 的映射。又设 w 是从 E 到 F 中把每个多项式 \(f\in E\) 对应于它在 \([0,1]\) 上的限制的映射。证明：u 是从 E 到 F 上的一个严格态射，但对每个 \(\varepsilon>0\) ，\(u+\varepsilon w\) 都不是从 E 到 F 中的严格态射。
\item
  \begin{enumerate}
  \def\labelenumii{\alph{enumii})}
  \tightlist
  \item
    设 E、F 是两个 Banach 空间，u 是从 E 到 F 中的一个严格态射，且 \(u^{-1}(0)\) 是有限维的。证明：对每个范数足够小的从 E 到 F 中的连续线性映射 w，\(v = u + w\) 都是从 E 到 F 中的一个严格态射，且 \(\dim v^{-1}(0) \leqslant \dim u^{-1}(0)\) 。（把 E 写成 \(u^{-1}(0)\) 与一个闭子空间的拓扑直和，并利用 IV, 第 66 页的 exerc. 18。）
  \end{enumerate}
\end{enumerate}

\begin{enumerate}
\def\labelenumi{\alph{enumi})}
\setcounter{enumi}{1}
\tightlist
\item
  设 E、F 是两个 Banach 空间，u 是从 E 到 F 中的一个连续线性映射，使得 \(u(\mathrm{E})\) 在 F 中有有限余维数。则 \(u(\mathrm{E})\) 是闭的，且 u 是一个严格态射 (I, 第 28 页, exerc. 4)。证明：对每个范数足够小的从 E 到 F 中的连续线性映射 w，\(v = u + w\) 都是从 E 到 F 中的一个严格态射，并且
\end{enumerate}

\[
\operatorname{codim} (v (\mathrm{E})) \leqslant \operatorname{codim} (u (\mathrm{E}))
\]

（考虑 \({}^t v = {}^t u + {}^t w\) ，并利用 a) 与 IV, 第 30 页, 推论 3）。

\begin{enumerate}
\def\labelenumi{\arabic{enumi})}
\setcounter{enumi}{20}
\tightlist
\item
  设 E 与 F 是两个 Banach 空间。我们称从 E 到 F 中的一个连续线性映射为 Fredholm 算子（或拟同构），如果 \(u^{-1}(0)\) 是有限维的且 \(u(\mathrm{E})\) 有有限余维数（这蕴含 \(u(\mathrm{E})\) 在 F 中是闭的且 u 是一个严格态射）；数 \(\operatorname{Ind}(u) = \operatorname{codim}(u(\mathrm{E})) - \dim(u^{-1}(0))\) 称为 u 的指标。
\end{enumerate}

\begin{enumerate}
\def\labelenumi{\alph{enumi})}
\item
  证明：\({}^t u: \mathrm{F}' \to \mathrm{E}'\) 也是一个 Fredholm 算子，并且 \(\operatorname{Ind}({}^t u) = -\operatorname{Ind}(u)\) 。
\item
  若 \(u: \mathrm{E} \to \mathrm{F}\) 与 \(v: \mathrm{F} \to \mathrm{G}\) 是两个 Fredholm 算子，则 \(v \circ u: \mathrm{E} \to \mathrm{G}\) 也是，且 \(\operatorname{Ind}(v \circ u) = \operatorname{Ind}(u) + \operatorname{Ind}(v)\) 。
\item
  若 \(w: E \to F\) 是一个有限秩或范数足够小的连续线性映射，则 \(u + w\) 是一个 Fredholm 算子，且 \(\operatorname{Ind}(u + w) = \operatorname{Ind}(u)\) （利用 IV, 第 65 页的 exerc. 17, c) 及 18, b)）。
\end{enumerate}

\begin{enumerate}
\def\labelenumi{\arabic{enumi})}
\setcounter{enumi}{21}
\tightlist
\item
  设 \(\mathbf{X}\) 是一个实 Banach 空间，\(\mathbf{E}\) 是 \(\mathbf{X}\) 的一个有限维子空间。
\end{enumerate}

\begin{enumerate}
\def\labelenumi{\alph{enumi})}
\item
  设 S 是 X 中的单位球面。证明：对每个 \(\varepsilon > 0\) ，存在有限多个线性无关的点 \(z_{i} \in S\) （ \(1 \leqslant i \leqslant r\) ），使得对每个 \(x \in S \cap E\) ，存在一个指标 \(i\) 使得 \(\| x - z_{i} \| \leqslant \varepsilon\) 。
\item
  设 \(z_{i}^{\prime}(1 \leqslant i \leqslant r)\) 是对偶 \(E^{\prime}\) （\(E\) 的对偶）中的点，满足 \(\| z_{i}^{\prime}\| \geqslant 1\) 对每个 \(i\) ，且 \(\langle z_i, z_j' \rangle = \delta_{ij}\) （Kronecker 符号），并设 \(F\) 是余维数为 \(r\) 的闭子空间（\(E\) 中的），它与 \(\mathbf{E}'\) 中由诸 \(z_{i}^{\prime}\) 生成的子空间正交。证明：对每个 \(x\in \mathrm{S}\cap \mathrm{E}\) 以及每个 \(y\in \mathrm{F}\) ，\(\| x + y\| \geqslant 1 - \varepsilon\) ，特别地 \(\mathrm{E}\cap \mathrm{F} = \{0\}\) ，从而和 \(\mathrm{E} + \mathrm{F}\) 是拓扑直和。
\item
  由 \(b)\) 推出：连续投影 \(P\) （从 \(\mathrm{E} + \mathrm{F}\) 到 \(\mathrm{E}\) 上的、对应于拓扑直和分解 \(\mathrm{E} \oplus \mathrm{F}\) 的）的范数 \(\| P \| < 1 / (1 - \varepsilon)\) 。
\end{enumerate}

\begin{enumerate}
\def\labelenumi{\arabic{enumi})}
\setcounter{enumi}{22}
\tightlist
\item
  设 \(X\) 是一个无限维实 Banach 空间。
\end{enumerate}

\begin{enumerate}
\def\labelenumi{\alph{enumi})}
\item
  假设对每个 \(\lambda > 0\) ，存在 X 的一个有限维子空间 \(\mathrm{E}_{\lambda}\) ，使得不存在 X 上的连续投影 \(P\) ，它以 \(\mathrm{E}_{\lambda}\) 为像且满足 \(\| P \| \leqslant \lambda\) 。证明：X 的每个有限余维数的闭子空间 Y 都具有与 X 相同的性质。
\item
  对 n 用归纳法证明：存在 X 的闭子空间（均有有限余维数）的一个递减序列 \((\mathrm{X}_{n})\) ，并且对每个 n，存在 X 的一个有限维子空间 \(E_{n}\) ，使得：\(1^{\circ}\) 和 \(E_{1} + E_{2} + \cdots + E_{n}\) 是直和，且存在一个连续投影 \(P_{n}\) （\(X_{n}\) 上的，范数 \(\leqslant 2\) ），其像为 \(E_{1} + E_{2} + \cdots + E_{n}\) ；\(2^{\circ}\) 空间 \(E_{n}\) 包含于 \((\mathrm{I} - P_{n-1})(\mathrm{X}_{n-1})\) 中；\(3^{\circ}\) 不存在 X 上以 \(E_{n}\) 为像且范数 \(\leqslant n + 2\) 的连续投影。
\item
  设 \(\bar{Z}\) 是 \(X\) 中由诸 \(E_{n}\) 的并生成的闭子空间。证明：\(Z\) 在 \(X\) 中不存在拓扑补（注意：若 \(Q\) 是 \(X\) 上的一个连续投影，以 \(\bar{Z}\) 为像，则 \((I - P_{n-1})P_nQ\) 将是 \(X\) 上以 \(E_{n}\) 为像的一个连续投影）。
\end{enumerate}
